\section{术语表}\etiquette[10]{I.sec10}

为方便读者, 这里汇集前面各节所用的一些术语的定义. 用 $C$ 表示一个范畴.

\subsection{笛卡儿, 余笛卡儿}\etiquette[10.1]{I.sec10.1}

$C$ 中的图表
\[
\xymatrix@C=3.25cm@R=1.5cm{
A\ar[r] \ar[d]& B \ar[d]\\
C\ar[r] & D }
\]
称为\emph{笛卡儿图表}, 如果它交换, 且从 $A$ 到纤维积 $C\times_D B$ 的典范态射是同构.
如果在 \sisi{$𝒜$}{$C$} 的反范畴中相应的图表是笛卡儿图表, 就称它为\emph{余笛卡儿图表}.

\subsection{满态射, 严格满态射}\etiquette[10.2]{I.sec10.2}

等等\pageoriginale. 参见 \hyperref[I.sec10.3]{10.3}.

\subsection{满态射族, 严格满态射族}\etiquette[10.3]{I.sec10.3}

等等: $C$ 中具有相同靶的箭头族
$$
f_i:A_i\longrightarrow B, i\in I{\quoi},
$$
称为\emph{满态射族}, 如果对 $C$ 的每个对象 $C$, 由 $f_i$ 诱导的映射
\hypertarget{sga4.exp.I.equation.10.3.1}{}\begin{equation}
\Hom_{C}(B,C)\longrightarrow \prod_i
\Hom_{C}(A_i,C),\tag{10.3.1}
\end{equation}
是单射.

所考虑的族称为\emph{严格满态射族}, 如果映射 \hyperlink{sga4.exp.I.equation.10.3.1}{(10.3.1)} 的像由满足下述条件的族 $(g_i)$ 组成: 对 $C$ 的每个对象 $R$, 每对指标 $i,j\in I$, 以及每对满足 $f_iu=f_jv$ 的箭头 $u:R\rightarrow A_i$, $v:R\rightarrow A_j$, 也有 $g_iu=g_jv$. 当 $i,j\in I$ 的纤维积 $A_i\times_BA_j$ 可表时, 这等价于说 \hyperlink{sga4.exp.I.equation.10.3.1}{(10.3.1)} 的像由满足下述条件的 $(g_i)$ 组成: 对每对指标 $i,j\in I$, 有 $g_i\pr_1=g_j\pr_2$, 其中 $\pr_1,\pr_2$ 是 $A_i\times_BA_j$ 的两个投影. 如果满足上述条件, 即它是严格满态射族且纤维积 $A_i\times_BA_j$ 可表, 则称族 $(f_i)_{i\in I}$ 为\emph{有效满态射族}.

族 $(f_i)_{i\in I}$ 称为\emph{泛满态射族} (分别称为\emph{泛有效满态射族}), 如果态射 $f_i$ 可作拉回 (\hyperref[I.sec10.3]{10.3}), 并且对每个箭头 $B'\rightarrow B$, 从族 $(f_i)_{i\in I}$ 经基变换 $B'\rightarrow B$ 得到的箭头族 $f'_i:A'_i=A_i\times_BB'\rightarrow B'$, $i\in I$, 是满态射族 (分别是有效满态射族).

\emph{泛严格满态射族}\pageoriginale 的概念将在 (II \hyperref[II.def2.5]{2.5}, \hyperref[II.prop2.6]{2.6}) 中定义. 注意, 当态射 $f_i$ 可作拉回时 (例如, 当 $C$ 中的纤维积可表时), 这个概念与泛有效满态射族的概念相同 (用 II \hyperref[II.cor2.4]{2.4}). 在一般情形, 所考虑的满性概念的六种变体之间有如下逻辑蕴涵:
\[
\xymatrix@C=-1cm@R=.5cm{
  &   & \txt<3cm>{泛有效满态射族} \ar@{=>}[lld]\ar@{=>}[rrd] &     & &\\
\txt<3cm>{有效满态射族} \ar@{=>}[rdd]       &   &    &     &  \txt<3cm>{泛严格满态射族} \ar@{=>}[rd] \ar@{=>}[ld]&\\
 &   &    &   \txt{严格\\ 满态射族}
         \ar@{=>}[lld]&     & \txt{泛 \\满态射族} \ar@{=>}[llld]\\
         & \txt{满态射族}&    &     &      &
}
\]
$C$ 的态射 $f:A\rightarrow B$ 称为\emph{满态射 (分别称为严格满态射、有效满态射、泛满态射、泛有效满态射、泛严格满态射)}, 如果仅由元素 $f$ 组成的态射族是满态射族 (分别为~\ldots).


% END INCLUDED SGA4_English_current_buildable_source_20260906/source/01/drafts/SGA4_I_LINES_7424_7512_English.texfrag


% BEGIN INCLUDED SGA4_English_current_buildable_source_20260906/source/01/drafts/SGA4_I_LINES_7513_7636_English.texfrag
\subsection{单态射族, 严格单态射族}\etiquette[10.4]{I.sec10.4}

等等. $C$ 中的箭头族 $f_i:A_i\rightarrow B$ 称为\emph{单态射族 (分别称为严格单态射族,\ldots)}, 如果把它看成反范畴 $C^{\circ}$ 中的箭头族时, 它是满态射族 (分别为严格满态射族,\ldots). $C$ 的箭头 $f:A\rightarrow B$ 称为\emph{单态射 (分别称为严格单态射,\ldots)}, 如果仅由 $f$ 组成的族是单态射族 (分别为\pageoriginale 严格单态射族,\ldots), 即如果把 $f$ 看成反范畴 $C^{\circ}$ 的箭头时, 它是满态射 (分别为严格满态射,\ldots).

\subsection{单态射, 严格单态射}\etiquette[10.5]{I.sec10.5}

等等. 参见 \hyperref[I.sec10.4]{10.4}.

\subsection{投影子, 投影子的像, 对象的直因子}\etiquette[10.6]{I.sec10.6}

设 $f:X\rightarrow X$ 为 $C$ 中一个对象的自同态. 若 $f^2=f$, 就称自同态 $f$ 为\emph{投影子}. 此时, 对 $(f,\id_X)$ 有可表的余核 $X'$ 当且仅当它有可表的核 $X''$; 在这种情形下, $C$ 中存在唯一的同构 $u:X'\rightarrow X''$ 使得 $iup=f$, 其中 $p:X\rightarrow X'$ 与 $i:X''\rightarrow X$ 是典范态射.\nde{核或余核的可表性, 也等价于存在态射 $i$, $p$ 和同构~$u$, 使得 $iup=f$ 且 $pi=u^{-1}$. 特别地, 每个函子都保持这一性质.}
此时通常将 $X'$ 与 $X''$ 等同, 并称之为\emph{投影子 $f$ 的像}; 若 $\Ker(f,\id_X)$ 可表, 即若 $\Coker(f,\id_X)$ 可表, 就说投影子 $f$ \emph{有像}. $X$ 的子对象 (分别为商) $Y$ 称为 $X$ 的\emph{直因子}, 如果存在 $X$ 上的投影子 $f$, 它经过 $Y$ 分解, 并且以 $Y$ 为它的像, 看成 $X$ 的子对象 (分别看成商). 有时也滥用语言, 把一个与 $X$ 上某投影子的像同构的 $C$ 的对象称为 $X$ 的\emph{直因子}.

\subsection{可作拉回}\etiquette[10.7]{I.sec10.7}

$C$ 的箭头 $f:A\rightarrow B$ 称为\emph{可作拉回的}, 如果对 $C$ 的每个箭头 $B'\rightarrow B$, 纤维积 $A\times_B B'$ 在 $C$ 中可表.

\subsection{商, 严格商}\etiquette[10.8]{I.sec10.8}

等等\pageoriginale.

设 $A$ 为 $C$ 的对象. 以 $A$ 为源的两个满态射 $f:A\rightarrow B$ 与 $f':A\rightarrow B'$ 称为\emph{等价的}, 如果存在同构 $u:B\xrightarrow{\sim}B'$ 使得 $uf=f'$. 这在以 $A$ 为源的满态射集上定义了一个等价关系, 它的等价类称为 $A$ 的\emph{商, 或商对象}. 对 $A$ 的每个商, 通常在相应的类中选取一个元素, 例如 $f:A\rightarrow B$, 并且常常 (滥用语言) 谈论 $A$ 的商 $B$ (或 $A$ 的商 $f:A\rightarrow B$). 商 $B$ 称为\emph{严格商} (分别称为\emph{有效商、泛商},\ldots), 如果态射 $f:A\rightarrow B$ 是严格满态射 (分别为有效满态射、泛满态射,\ldots) (\hyperref[I.sec10.3]{10.3}).

\subsection{等价关系}\etiquette[10.9]{I.sec10.9}

设 $F$ 与 $G$ 是 $C$ 上的两个集合预层. 图表 $p_1$, $p_2:F\rightrightarrows G$ 称为 $G$ 上的\emph{等价关系}, 如果对 $C$ 的每个对象 $A$, 映射
$$
(p_1(A),p_2(A)):F(A)\longrightarrow G(A)\times G(A)
$$
诱导 $F(A)$ 到集合 $G(A)$ 上某个等价关系的图的双射.

$C$ 中的图表 $p_1$, $p_2:B\rightrightarrows C$ 称为 $C$ 上的等价关系, 如果相应的预层图表是等价关系. 当 $C$ 中的有限积与纤维积可表时, 与有限积和纤维积交换的函子 $\Phi:C\rightarrow C'$ 把 $C$ 上的等价关系映成 $\Phi(C)$ 上的等价关系.

设\pageoriginale $\pi:C\rightarrow D$ 是 $C$ 的态射, 且纤维积 $C\times_D C$ 可表. 则图表 $\pr_1$, $\pr_2:C\times_D C\rightrightarrows C$ 是~$C$ 上的等价关系.

\subsection{有效等价关系, 泛有效等价关系}\etiquette[10.10]{I.sec10.10}

$C$ 的对象 $A$ 上的等价关系 $p_1$, $p_2:R\rightrightarrows A$ 称为\emph{有效等价关系}, 如果存在态射 $\pi:A\rightarrow B$, 使方块
\[
\xymatrix@C=3.75cm@R=1.5cm{
R\ar[r]^{p_2} \ar[d]_{p_1}& A \ar[d]^{\pi}\\
A\ar[r] & B }
\]
既是笛卡儿的, 又是余笛卡儿的. (态射 $\pi$ 是对 $(p_1,p_2)$ 的余核, 因而在唯一同构的意义下确定.) 此时态射 $\pi$ 是有效满态射; 如果它是泛有效满态射 (\hyperref[I.sec10.3]{10.3}), 就说此\emph{等价关系}是\emph{泛有效的}.

\subsection{子对象, 严格子对象}\etiquette[10.11]{I.sec10.11}

等等. 这些是商、严格商等概念 (\hyperref[I.sec10.8]{10.8}) 的对偶概念.

\pdfbookmark[1]{参考文献}{sga4.exp.I.bibliography.main}
\begingroup
\makeatletter
\def\Hy@writebookmark#1#2#3#4#5{}
\makeatother
\begin{thebibliography}{1}
\bibitem[1]{sga4.exp.I.bib.I.1} \sisi{B. Mitchell, Theory of Categories, Academic Press
(1965).}{\textsc{B. Mitchell}, ``Theory of Categories,'' Academic
Press (1965).}
\end{thebibliography}
\endgroup


\newpage

\footnotetext[1]{这里经 N. \textsc{Bourbaki} 许可转载他的秘密文稿. 本文中的文献引用均指向他的博学著作.}

% END INCLUDED SGA4_English_current_buildable_source_20260906/source/01/drafts/SGA4_I_LINES_7513_7636_English.texfrag


% BEGIN INCLUDED SGA4_English_current_buildable_source_20260906/source/01/drafts/SGA4_I_LINES_7637_7820_English.texfrag
\pdfbookmark[1]{II. 附录: 宇宙 (N. Bourbaki 著)}{sga4.exp.I.appendix.II}
\begingroup
\makeatletter
\def\Hy@writebookmark#1#2#3#4#5{}
\makeatother
\section*{II. 附录: 宇宙 (N. \textsc{Bourbaki} 著($^*$))}\pageoriginale
\endgroup
\markboth{II. 附录: 宇宙}{II. 附录: 宇宙}

\hypersetup{bookmarksdepth=2}
\makeatletter
\let\sgaFourSavedToclevelSection\toclevel@section
\def\toclevel@section{2}
\makeatother

\setcounter{section}{0}
\section{宇宙的定义与初步性质}%
\hypertarget{sga4.exp.I.appendix.II.sec.1}{}

\hypertarget{sga4.exp.I.appendix.II.sec.1.definition.1}{}
\begin{enonce*}{定义 1}
集合 $U$ 称为宇宙, 如果它满足以下条件:
\begin{itemize}
\hypertarget{sga4.exp.I.appendix.II.sec.1.definition.1.condition.U.I}{}
\item[(U.I)] 若 $x\in U$ 且 $y\in x$, 则 $y\in U$;

\hypertarget{sga4.exp.I.appendix.II.sec.1.definition.1.condition.U.II}{}
\item[(U.II)] 若 $x$, $y\in U$, 则 $\{x,y\}\in U$;

\hypertarget{sga4.exp.I.appendix.II.sec.1.definition.1.condition.U.III}{}
\item[(U.III)] 若 $x\in U$, 则 $\mathcal{P}(x)\in U$;

\hypertarget{sga4.exp.I.appendix.II.sec.1.definition.1.condition.U.IV}{}
\item[(U.IV)] 若 $(x_{\alpha})_{\alpha\in I}$ 是 $U$ 的元素族, 且 $I\in U$, 则并集 $\bigcup_{\alpha\in I}x_{\alpha}$ 属于 $U$;

\hypertarget{sga4.exp.I.appendix.II.sec.1.definition.1.condition.U.question}{}
\item[(U.?)] 若 $x$, $y\in U$, 则有序对 $(x,y)\in U$.
\end{itemize}
\end{enonce*}
\noindent 注意. 我想, 既已决定在今后的版本中用 $(x,y)=\{\{x,y\}, \{x\}\}$ 按 \sisi{Kuratowaki}{Kuratowski} 的方式定义有序对, 条件 (U.?) 就不必列出了, 因为它可由 \hyperlink{sga4.exp.I.appendix.II.sec.1.definition.1.condition.U.II}{(U.II)} 推出.

\hypertarget{sga4.exp.I.appendix.II.sec.1.examples}{}
\begin{enonce*}[remark]{例}
\begin{enumerate}
\renewcommand{\labelenumi}{\theenumi)}
\hypertarget{sga4.exp.I.appendix.II.sec.1.example.1}{}
\item \emph{空}集是宇宙, 记为 $U_\circ$.

\hypertarget{sga4.exp.I.appendix.II.sec.1.example.2}{}
\item 考虑用四个符号 $''\{'',''\}'','',''$ 和 $''∅''$ 构成的非空有限字 (参见 Alg. I). 对这种字的长度 $n$ 作归纳\pageoriginale, 定义\emph{有意义的}字:
\begin{enumerate}
\renewcommand{\theenumii}{\alph{enumii}}
\renewcommand{\labelenumii}{(\theenumii)}
\item 当 $n=1$ 时, 只有字 $∅$ 是有意义的;

\item 长度为 $n$ 的字 $A$ 有意义的充分必要条件是, 存在 $p$ 个长度 $<n$ 的互异的有意义的字 $A_1,\ldots,A_p$ $(p\geq 1)$, 使得
$$
A=\{A_1,A_2,\ldots,A_p\}{\quoi}.
$$
\end{enumerate}
\end{enumerate}
\end{enonce*}
\noindent 例如, $∅$, $\{∅\}$, $\{\{∅\},\{\{∅\}\},\{∅,\{∅\}\}\}$ 都是有意义的字. 对长度作归纳, 显然每个有意义的字都表示集合论中的一个项; 例如, 如果 $A=\{A_1,A_2,\ldots,A_p\}$, 则 $A$ 表示以 $A_1$, $A_2,\ldots,$ $A_p$ 为元素的集合; 这些项显然都是\emph{有限集}. 设 $U_1$ 为这样得到的集合所组成的集合. 容易验证 $U_1$ 满足\hyperlink{sga4.exp.I.appendix.II.sec.1.definition.1}{定义 1} 的条件 \hyperlink{sga4.exp.I.appendix.II.sec.1.definition.1.condition.U.I}{(U.I)}, \hyperlink{sga4.exp.I.appendix.II.sec.1.definition.1.condition.U.II}{(U.II)}, \hyperlink{sga4.exp.I.appendix.II.sec.1.definition.1.condition.U.III}{(U.III)} 和 \hyperlink{sga4.exp.I.appendix.II.sec.1.definition.1.condition.U.IV}{(U.IV)} (但不满足那个愚蠢的条件 (U.?); 只要肯去掉有序对的管子表示, 这就不成问题). 因此 $U_1$ 是一个\emph{宇宙}. 注意 $U_1$ 的元素都是有限集, 而 $U_1$ 是可数的.

在以下各陈述中, $U$ \emph{表示一个宇宙}.

\hypertarget{sga4.exp.I.appendix.II.sec.1.proposition.1}{}
\begin{enonce*}{命题 1}
若 $X\in U$ 且 $y\subset x$, 则 $y\in U$.
\end{enonce*}
事实上, 由 \hyperlink{sga4.exp.I.appendix.II.sec.1.definition.1.condition.U.III}{(U.III)} 有 $\mathcal{P}(x)\in U$, 因此 $y\in\mathcal{P}(x)$, 再由 \hyperlink{sga4.exp.I.appendix.II.sec.1.definition.1.condition.U.I}{(U.I)} 得 $y\in U$.

\hypertarget{sga4.exp.I.appendix.II.sec.1.corollary.to.proposition.1}{}
\begin{enonce*}{推论}
若 $x\in U$, 则 $x$ 的每个商集 $y$ 都是 $U$ 的元素.
\end{enonce*}
事实上, $y$ 是 $\mathcal{P}(x)$ 的子集. 因而由 \hyperlink{sga4.exp.I.appendix.II.sec.1.definition.1.condition.U.III}{(U.III)} 和\hyperlink{sga4.exp.I.appendix.II.sec.1.proposition.1}{命题 1} 得 $y\in U$.

\hypertarget{sga4.exp.I.appendix.II.sec.1.proposition.2}{}
\begin{enonce*}{命题 2}
若\pageoriginale $x\in U$, 则 $\{x\}\in U$.
\end{enonce*}
在 \hyperlink{sga4.exp.I.appendix.II.sec.1.definition.1.condition.U.II}{(U.II)} 中取 $y=x$, 即得此结论.

\hypertarget{sga4.exp.I.appendix.II.sec.1.proposition.3}{}
\begin{enonce*}{命题 3}
由 $U$ 的元素组成的每个有序对、有序三元组及有序四元组都是 $U$ 的元素.
\end{enonce*}
由上述注意或 \hyperlink{sga4.exp.I.appendix.II.sec.1.definition.1.condition.U.question}{(U.?)} 可知这对有序对成立. 因为 $(x,y,z)=((x,y),z)$, 有序三元组的情形随之得证; 再由 $(x,y,z,t)=((x,y,z),t)$ 得有序四元组的情形.

\hypertarget{sga4.exp.I.appendix.II.sec.1.proposition.4}{}
\begin{enonce*}{命题 4}
若 $X$, $Y\in U$, 则 $X\times Y\in U$.
\end{enonce*}
事实上, 对 $x\in X$, $\{x\}\times Y$ 是族 $\{(x,y)\}_{y\in Y}$ 的并集; 由 \hyperlink{sga4.exp.I.appendix.II.sec.1.definition.1.condition.U.I}{(U.I)}, \hyperlink{sga4.exp.I.appendix.II.sec.1.definition.1.condition.U.II}{(U.II)} 和\hyperlink{sga4.exp.I.appendix.II.sec.1.proposition.3}{命题 3} 有 $\{(x,y)\}\in U$, 因此由 (U.IV) 有 $\{x\}\times Y\in U$. 最后, $X\times Y$ 是族 $(\{x\}\times Y)_{x\in X}$ 的并集; 再由 \hyperlink{sga4.exp.I.appendix.II.sec.1.definition.1.condition.U.IV}{(U.IV)} 可知它是 $U$ 的元素.

\hypertarget{sga4.exp.I.appendix.II.sec.1.corollary.to.proposition.4.1}{}
\begin{enonce*}{推论 1}
若 $X$, $Y$, $Z,\ldots$ 是 $U$ 的元素, 则以 $X$, $Y$, $Z,\ldots$ 为基集构造的阶梯中的一切集合都是 $U$ 的元素 {\rm(参见第 IV 章, \S)}.
\end{enonce*}

依次应用\hyperlink{sga4.exp.I.appendix.II.sec.1.proposition.4}{命题 4} 和 (U.III), 即得此结论.

\hypertarget{sga4.exp.I.appendix.II.sec.1.corollary.to.proposition.4.2}{}
\begin{enonce*}{推论 2}
若 $(x_{\alpha})_{\alpha \in I}$ 为 $U$ 的元素族, 且 $I\in U$, 则和集 $\Sigma_{\alpha \in I} x_{\alpha}$ 是 $U$ 的元素.
\end{enonce*}
事实上, 这个和集是积 $(\bigcup_{\alpha\in I} x_{\alpha})\times I$ 的子集 (第 II 章), 而该积的两个因子都是 $U$ 的元素 (由 (U.IV) 与假设). 再应用命题 \hyperlink{sga4.exp.I.appendix.II.sec.1.proposition.4}{4} 和 \hyperlink{sga4.exp.I.appendix.II.sec.1.proposition.1}{1} 即可.

\hypertarget{sga4.exp.I.appendix.II.sec.1.proposition.5}{}
\begin{enonce*}{命题 5}
若\pageoriginale $X$ 与 $Y$ 是 $U$ 的元素, 则 $X$ 与 $Y$ 之间的每个对应 (特别是从 $X$ 到 $Y$ 的每个映射) 都是 $U$ 的元素.
\end{enonce*}
事实上, 这样的对应 $C$ 是有序三元组 $(X,Y,\Gamma)$, 其中 $\Gamma$ 是 $X\times Y$ 的子集 ($C$ 的图) (第 II 章, \S). 由命题 \hyperlink{sga4.exp.I.appendix.II.sec.1.proposition.4}{4} 和 1 得 $\Gamma\in U$. 因此由\hyperlink{sga4.exp.I.appendix.II.sec.1.proposition.3}{命题 3} 得 $C\in U$.

\hypertarget{sga4.exp.I.appendix.II.sec.1.proposition.6}{}
\begin{enonce*}{命题 6}
若 $X$, $Y\in U$, 则由 $X$ 与 $Y$ 之间的对应 (特别是从 $X$ 到 $Y$ 的映射) 组成的每个集合 $Z$ 都是 $U$ 的元素.
\end{enonce*}
事实上, $Z$ 是 $\{X\}\times \{Y\}\times\mathcal{P}(X\times Y)$ 的子集. 而由\hyperlink{sga4.exp.I.appendix.II.sec.1.proposition.2}{命题 2} 和\hyperlink{sga4.exp.I.appendix.II.sec.1.corollary.to.proposition.4.1}{命题 4 的推论}可知此积是 $U$ 的元素. 因此由\hyperlink{sga4.exp.I.appendix.II.sec.1.proposition.1}{命题 1} 得 $Z\in U$.

\hypertarget{sga4.exp.I.appendix.II.sec.1.corollary.to.proposition.6}{}
\begin{enonce*}{推论}
若 $(x_{\alpha})_{\alpha \in I}$ 为 $U$ 的元素族, 且 $I\in U$, 则 $\Pi_{\alpha\in I}x_{\alpha} \in U$.
\end{enonce*}
事实上, 这个积是由从 $I$ 到 $\bigcup_{\alpha\in I}x_{\alpha}$ 的映射组成的集合, 而由 \hyperlink{sga4.exp.I.appendix.II.sec.1.definition.1.condition.U.IV}{(U.IV)} 可知此并集是 $U$ 的元素.

\hypertarget{sga4.exp.I.appendix.II.sec.1.proposition.7}{}
\begin{enonce*}{命题 7}
若 $X$ 是 $U$ 的子集, 其势不超过 $U$ 的某个元素的势, 则 $X$ 是 $U$ 的元素.
\end{enonce*}
设 $I$ 为 $U$ 的元素, 使得 $\card(X)\leq \card(I)$. 存在从 $I$ 的某个子集 $I'$ 到 $X$ 的满射 $i\mapsto x_i$. 此时 $X$ 是族 $(\{x_i\})_{i\in I'}$ 的并集; 由\hyperlink{sga4.exp.I.appendix.II.sec.1.proposition.1}{命题 1}、\hyperlink{sga4.exp.I.appendix.II.sec.1.proposition.2}{命题 2} 和 \hyperlink{sga4.exp.I.appendix.II.sec.1.definition.1.condition.U.IV}{(U.IV)}, 这个并集是 $U$ 的元素.

\hypertarget{sga4.exp.I.appendix.II.sec.1.corollary.to.proposition.7}{}
\begin{enonce*}{推论}
若\pageoriginale $U$ 非空, 则 $U$ 的每个有限子集都是 $U$ 的元素, 且 $U$ 含有有限势任意大的元素.
\end{enonce*}
另一方面, 若 $U=∅$, 则 $∅$ 是 $∅$ 的有限子集, 却不是 $∅$ 的元素.

事实上, 若 $U$ 非空, 则由\hyperlink{sga4.exp.I.appendix.II.sec.1.proposition.2}{命题 2} 可知有一个单元素集 $x_{\circ}$ 属于 $U$. 对 $n$ 归纳定义 $x_{n+1}=\mathcal{P}(x_n)$. 由 \hyperlink{sga4.exp.I.appendix.II.sec.1.definition.1.condition.U.III}{(U.III)} 有 $x_n\in U$, 且 $\card(x_{n+1})=2^{\card(x_n)}$, 因此 $U$ 含有有限势任意大的元素.

\begin{enonce*}[remark]{注}
由\hyperlink{sga4.exp.I.appendix.II.sec.1.proposition.2}{命题 2} 及\hyperlink{sga4.exp.I.appendix.II.sec.1.corollary.to.proposition.7}{命题 7 的推论}可知, 每个非空宇宙 $U$ 都包含\hyperlink{sga4.exp.I.appendix.II.sec.1.example.2}{例 2} 中的宇宙 $U_1$; 事实上, 由\hyperlink{sga4.exp.I.appendix.II.sec.1.proposition.1}{命题 1} 有 $∅\in U$. 因而 $U_1$ 是所有非空宇宙的交集. 更一般地:
\end{enonce*}

\hypertarget{sga4.exp.I.appendix.II.sec.1.proposition.8}{}
\begin{enonce*}{命题 8}
若 $(U_{\lambda})_{\lambda\in L}$ 是非空宇宙族, 则 $U=\bigcap_{\lambda\in L}U_{\lambda}$ 是宇宙.
\end{enonce*}
这由\hyperlink{sga4.exp.I.appendix.II.sec.1.definition.1}{定义 1} 立即得出.

% END INCLUDED SGA4_English_current_buildable_source_20260906/source/01/drafts/SGA4_I_LINES_7637_7820_English.texfrag


% BEGIN INCLUDED SGA4_English_current_buildable_source_20260906/source/01/drafts/SGA4_I_LINES_7821_7923_English.texfrag
\section{宇宙与结构种}%
\hypertarget{sga4.exp.I.appendix.II.sec.2}{}

设 $\mathcal{E}$ 为一个结构种; 为确定概念并简化叙述, 假定每个 $\mathcal{E}$ 种结构都定义在一个基集上. 设 $(X,S)$ 为一个 $\mathcal{E}$ 种结构 (其中 $X$ 为基集, $S$ 为结构), 并设 $U$ 为一个\emph{宇宙}. 若 $X\in U$, 则 $X$ 上结构 $S$ 的一切组成对象都是 $U$ 的元素 (由\hyperlink{sga4.exp.I.appendix.II.sec.1}{第 1 节}\hyperlink{sga4.exp.I.appendix.II.sec.1.corollary.to.proposition.4.1}{命题 4 的推论 1} 及 \hyperlink{sga4.exp.I.appendix.II.sec.1.definition.1.condition.U.I}{(U.I)}), 因此结构 $(X,S)$ 是 $U$ 的元素.

现在\pageoriginale 假定 $\mathcal{E}$ 为带有\emph{态射}的结构种. 若 $X$ 与 $X'$ 为 $U$ 的元素, 且赋予了 $\mathcal{E}$ 种结构, 则从 $X$ 到 $X'$ 的态射集仍是 $U$ 的元素 (\hyperlink{sga4.exp.I.appendix.II.sec.1}{第 1 节}, \hyperlink{sga4.exp.I.appendix.II.sec.1.proposition.6}{命题 6}).

现在考虑 \hyperlink{sga4.exp.I.appendix.II.sec.1}{\S 1}, 第 2 节, 例 c) 中定义的\emph{范畴} $(U\text{-}\mathcal{E})$. 因为 $(U\text{-}\mathcal{E})$ 的对象和箭头都是 $U$ 的元素, $(U\text{-}\mathcal{E})$ 是 $U$ 的\emph{一对子集}, 配有四个映射; 一般而言, 它不是 $U$ 的元素. 记号 $X\in (U\text{-}\mathcal{E})$ 表示 $X$ 是赋有 $\mathcal{E}$ 种结构的 $U$ 的元素.

范畴 $(U\text{-}\mathcal{E})$ 对许多运算\emph{稳定}:

\begin{enumerate}
\renewcommand{\theenumi}{\alph{enumi}}
\renewcommand{\labelenumi}{\theenumi)}
\item 若 $X\in(U\text{-}\mathcal{E})$, 且 $X'$ 是 $X$ 的一个容许\emph{诱导结构}的子集, 则 $X'\in(U\text{-}\mathcal{E})$. 这由\hyperlink{sga4.exp.I.appendix.II.sec.1}{第 1 节}\hyperlink{sga4.exp.I.appendix.II.sec.1.proposition.1}{命题 1} 得出.

\item 若 $X\in (U\text{-}\mathcal{E})$, 且 $X''$ 是 $X$ 的一个容许\emph{商结构}的商集, 则 $X''\in(U\text{-}\mathcal{E})$ (\hyperlink{sga4.exp.I.appendix.II.sec.1}{第 1 节}\hyperlink{sga4.exp.I.appendix.II.sec.1.corollary.to.proposition.1}{命题 1 的推论}).

\item 若 $X$, $Y\in (U\text{-}\mathcal{E})$, 且 $X\times Y$ 容许\emph{积结构}, 则 $X\times Y\in(U\text{-}\mathcal{E})$ (\hyperlink{sga4.exp.I.appendix.II.sec.1}{第 1 节}, \hyperlink{sga4.exp.I.appendix.II.sec.1.proposition.4}{命题 4}). 更一般地, 若 $(X_i)_{i\in I}$ 为 $(U\text{-}\mathcal{E})$ 的元素族, \emph{若} $I\in U$, 且 $\Pi_{i\in I}X_i$ 容许\emph{积结构}, 则 $\Pi_{i\in I}X_i\in (U\text{-}\mathcal{E})$ (\hyperlink{sga4.exp.I.appendix.II.sec.1.corollary.to.proposition.6}{命题 6 的推论}). 关于\emph{和结构}也有类似的断言 (\hyperlink{sga4.exp.I.appendix.II.sec.1.corollary.to.proposition.4.2}{命题 4 的推论 2}).

\item 设 $I$ 为预序集, 并设 $(X_i, f_{ij})_{i,j\in I}$ 为赋有 $\mathcal{E}$ 种结构及态射的集合的\emph{投射系} (分别为\emph{归纳系}); 设 $L$ 为第 III 章意义下此系的极限. 若对每个 $i$ 有 $X_i\in U$, 若 $I\in U$, 且 $L$ 容许投射极限结构 (分别为归纳极限结构), 则 $L\in (U\text{-}\mathcal{E})$: 事实上, $L$ 是 $\Pi_{i\in I}X_i$ 的子集\pageoriginale (分别为 $\Sigma_{i\in I}X_i$ 的商集), 因而由\hyperlink{sga4.exp.I.appendix.II.sec.1}{第 1 节}可知它是 $U$ 的元素.
\end{enumerate}

再举几个较具体的例子:
\begin{itemize}
\hypertarget{sga4.exp.I.appendix.II.sec.2.example.star.1}{}
\item[* 1)] 设 $X$, $Y\in (U\text{-}\Top)$ 为两个\emph{局部紧}空间. 则从 $X$ 到 $Y$ 的连续映射集 $\mathscr{C}(X,Y)$ 赋予紧收敛拓扑 (General Topology $X$) 后, 是 $(U\text{-}\Top)$ 的元素. 事实上, 由\hyperlink{sga4.exp.I.appendix.II.sec.1}{第 1 节}\hyperlink{sga4.exp.I.appendix.II.sec.1.proposition.6}{命题 6} 得 $\mathscr{C}(X,Y)\in U$.${}_{\mbox{*}}$

\hypertarget{sga4.exp.I.appendix.II.sec.2.example.star.2}{}
\item[* 2)] 设 $X$, $Y\in (U\text{-}\Ab)$. 则由\hyperlink{sga4.exp.I.appendix.II.sec.1}{第 1 节}\hyperlink{sga4.exp.I.appendix.II.sec.1.proposition.6}{命题 6} 可知群 $\Hom_{\ZZ}(X,Y)$ 是 $(U\text{-}\Ab)$ 的元素. 此外, 若 $\ZZ\in U$, 则张量积 $X\otimes_{\ZZ}Y$ 是 $(U\text{-}\Ab)$ 的元素; 事实上, 此张量积是 $\ZZ^{X\times Y}$ 的商, 后者是 $\ZZ^{X\times Y}$ 的子集, 再由\hyperlink{sga4.exp.I.appendix.II.sec.1}{第 1 节}命题 \hyperlink{sga4.exp.I.appendix.II.sec.1.proposition.4}{4} 和 \hyperlink{sga4.exp.I.appendix.II.sec.1.proposition.6}{6} 可知它是 $U$ 的元素.${}_{\mbox{*}}$

\hypertarget{sga4.exp.I.appendix.II.sec.2.example.star.3}{}
\item[* 3)] 设 $X\in (U\text{-}\Unif)$ 为一致空间. 则其\emph{完备化} $\widehat{X}$ 是 $(U\text{-}\Unif)$ 的元素. 事实上, $\widehat{X}$ 是 $X$ 上 Cauchy 滤子的等价类集; 而 $X$ 上的滤子是 $\mathcal{P}(\mathcal{P}(X))$ 的元素, 因此滤子的一个等价类是 $\mathcal{P}(\mathcal{P}(\mathcal{P}(X)))$ 的元素, 所以 $\widehat{X}$ 是 $\mathcal{P}(\mathcal{P}(\mathcal{P}(\mathcal{P}(X))))$ 的元素, 从而由 \hyperlink{sga4.exp.I.appendix.II.sec.1.definition.1.condition.U.III}{(U.III)} 及 \hyperlink{sga4.exp.I.appendix.II.sec.1.definition.1.condition.U.I}{(U.I)} 可知它是 $U$ 的元素.${}_{\mbox{*}}$
\end{itemize}

注意. 寓意: Bourbaki 必须用心把其构造适当地「典范化」. 例如, 在添入一个元素 $\infty$ 的构造中 (Alexandroff 紧化、射影域), 取 $\infty =∅$\pageoriginale (因为 $∅$ 是每个非空宇宙的元素), 并构造 $\{∅\}$ 与给定集合的不交并, 会比较有利. 当然, 用于构造不交并的二元素集也应「典范化」; 因而在我看来, $\{∅, \{∅\}\}$ 是一个很好的候选, 因为它属于每个非空宇宙.

% END INCLUDED SGA4_English_current_buildable_source_20260906/source/01/drafts/SGA4_I_LINES_7821_7923_English.texfrag


% BEGIN INCLUDED SGA4_English_current_buildable_source_20260906/source/01/drafts/SGA4_I_LINES_7924_7993_English.texfrag
\section{宇宙与范畴}%
\hypertarget{sga4.exp.I.appendix.II.sec.3}{}

设 $U$ 为宇宙, $C$ 为范畴. 如果 $\Ob(C)\in U$ 且 $%%
\mathrm{F}{}l\,(C)\in U$, 就记 $C\in U$. 这个记号是有根据的, 因为 $C$ 是由 $\Ob(C)$, $\mathrm{F}{}l\,(C)$ 及四个结构映射组成的六元组; 因此, 若 $\Ob(C)\in U$ 且 $\mathrm{F}{}l\,(C)\in U$, 则这四个映射都是 $U$ 的元素 (\hyperlink{sga4.exp.I.appendix.II.sec.1}{第 1 节}, \hyperlink{sga4.exp.I.appendix.II.sec.1.proposition.4}{命题 4} 的\hyperlink{sga4.exp.I.appendix.II.sec.1.corollary.to.proposition.4.1}{推论 1}), 从而六元组 $C$ 也是 $U$ 的元素 (\hyperlink{sga4.exp.I.appendix.II.sec.1}{第 1 节}, \hyperlink{sga4.exp.I.appendix.II.sec.1.proposition.3}{命题 3}, \textit{cum grano salis}). 此外, 如果考虑范畴的结构种「$\cat$」, 这也是\hyperlink{sga4.exp.I.appendix.II.sec.2}{第 2 节} 的一个特例. 按\hyperlink{sga4.exp.I.appendix.II.sec.2}{第 2 节}的记号, 关系 $C\in U$ 也记为 $C\in(U\text{-}\cat)$.

应当注意, 像 $(U\text{-}\Ens)$, $(U\text{-}\Ord)$, $(U\text{-}\Top)$ 或 $(U\text{-}\Gr)$ 这样的范畴, 一般\emph{不是} $U$ 的元素.

设 $C,D$ 为两个范畴, $U$ 为宇宙, 且 $C,D\in U$. 若 $C'$ 为 $C$ 的\emph{子范畴}, 则 $C'\in U$; \emph{积范畴} $C\times D$、$C$ 与 $D$ 的和范畴, 以及反范畴 $C^\circ$ 与 $D^\circ$ 也都是 $U$ 的元素 (参见\hyperlink{sga4.exp.I.appendix.II.sec.2}{第 2 节}). \emph{函子范畴} $E=\Hhom(C,D)$ 同样是 $U$ 的元素: 事实上, $\Ob(E)$ 是由映射对 $\Ob(C)\rightarrow\Ob(D)$, $F\ell\,(C)\rightarrow F\ell\,(D)$ 组成的集合, 因而由\hyperlink{sga4.exp.I.appendix.II.sec.1}{第 1 节}得 $\Ob(E)\in U$;\pageoriginale 至于 $F\ell\,(E)$, 它是自然变换集, 即映射 $\Ob(C)\rightarrow F\ell\,(D)$ 的集合, 因而由\hyperlink{sga4.exp.I.appendix.II.sec.1}{第 1 节}\hyperlink{sga4.exp.I.appendix.II.sec.1.proposition.6}{命题 6} 得 $F\ell\,(E)\in U$.

\hypertarget{sga4.exp.I.appendix.II.sec.3.proposition.9}{}
\begin{enonce*}{命题 9}
设 $\mathcal{E}$ 为带有态射的结构种, 并设 $U,U'$ 为两个宇宙, 且 $U'\subset U$. 则 $(U'\text{-}\mathcal{E})$ 是 $(U\text{-}\mathcal{E})$ 的满子范畴.
\end{enonce*}

事实上, 若 $x$ 与 $y$ 是 $(U'\text{-}\mathcal{E})$ 的对象, 则按定义它们是 $(U\text{-}\mathcal{E})$ 的对象. 至于 $\Hom(x,y)$, 在两个范畴中它都是同一个集合 (参见 \hyperlink{sga4.exp.I.appendix.II.sec.1}{\S 1}, 第 2 节, 例 c)).

注意. 应当注意, $U$ 与 $U'$ 为宇宙的假设并不必要, 所以\hyperlink{sga4.exp.I.appendix.II.sec.3.proposition.9}{命题 9} 放在 \hyperlink{sga4.exp.I.appendix.II.sec.1}{\S 1}, 第 4 节中会更合适. 不过部落要求把它放在这里.

\begin{enonce*}[remark]{注}
结构种 $\mathcal{E}$ 的公理可能蕴涵: 赋有 $\mathcal{E}$ 种结构的集合的\emph{基数}, 以某个固定基数, 例如 $\sheaf{\sheaf{c}}$, 为\emph{界} (* 例如有限群、有限生成群、固定环 $A$ 上的有限生成模或 $A$ 上的有限生成代数的结构种 ${}_{\mbox{*}}$). 此时假设存在 $U'$ 的元素 $z$ 使得 $\sheaf{\sheaf{c}}\leq\card(z)$. 则对每个赋有 $\mathcal{E}$ 种结构的集合 $x$, 都存在与 $x$ 等势的 $U'$ 的元素 $x'$, 例如 $z$ 的一个子集; 给 $x'$ 赋予从 $x$ 的结构迁移而来的结构, 就得到 $(U'\text{-}\mathcal{E})$ 的元素. 由此可知, \emph{从 $(U'\text{-}\mathcal{E})$ 到 $(U\text{-}\mathcal{E})$ 的包含函子此时本质满射 {\rm(\hyperlink{sga4.exp.I.appendix.II.sec.4}{\S 4}, 第 2 节, 定义 2)}, 因而是范畴等价} (\hyperlink{sga4.exp.I.appendix.II.sec.4}{\S 4}, 第 3 节, 定理 1).
\end{enonce*}

% END INCLUDED SGA4_English_current_buildable_source_20260906/source/01/drafts/SGA4_I_LINES_7924_7993_English.texfrag


% BEGIN INCLUDED SGA4_English_current_buildable_source_20260906/source/01/drafts/SGA4_I_LINES_7994_8025_English.texfrag
\section{宇宙公理}%
\hypertarget{sga4.exp.I.appendix.II.sec.4}{}\pageoriginale

\hyperlink{sga4.exp.I.appendix.II.sec.2}{第 2 节}中那些十分可喜的稳定性结果, 只有在能够应用于\hyperlink{sga4.exp.I.appendix.II.sec.1}{第 1 节}所描述的两个小宇宙 $U_\circ$ 与 $U_1$ 之外的对象时才有意义. 因此, 我们将在集合论的公理中添入以下公理:

\hypertarget{sga4.exp.I.appendix.II.sec.4.axiom.A.6}{}
(A.6) \emph{对每个集合 $x$, 都存在宇宙 $U$ 使得 $x\in U$.}

这个公理蕴涵: 若 $(x_\alpha)_{\alpha\in I}$ 是集合\emph{族}, 则存在宇宙 $U$, 使对每个 $\alpha\in I$ 都有 $x_\alpha\in U$: 事实上, 只须对 $x=\bigcup_{\alpha\in I}x_\alpha$ 应用 \hyperlink{sga4.exp.I.appendix.II.sec.4.axiom.A.6}{(A.6)}, 再应用\hyperlink{sga4.exp.I.appendix.II.sec.1}{第 1 节}\hyperlink{sga4.exp.I.appendix.II.sec.1.proposition.1}{命题 1} 即可.

特别地, 给定\emph{范畴} $C$, 存在宇宙 $U$, 使得在\hyperlink{sga4.exp.I.appendix.II.sec.3}{第 3 节}的意义下 $C\in U$: 对族 $(\Ob(C),F\ell\,(C))$ 应用上述结果即可. 这适用于形如 $(V\text{-}\mathcal{E})$ 的范畴, 其中 $\mathcal{E}$ 为带有态射的结构种, $V$ 为集合, 例如宇宙; 一般有 $V\neq U$.

例如, 若 $V$ 为宇宙, 则 $\Ob(V\text{-}\Ens)=V$ 且 $F\ell\,(V\text{-}\Ens)\subset V$ (\hyperlink{sga4.exp.I.appendix.II.sec.1}{第 1 节}, \hyperlink{sga4.exp.I.appendix.II.sec.1.proposition.5}{命题 5}). 因而使得 $(V\text{-}\Ens)\in U$ 的宇宙 $U$, 正是使得 $V\in U$ 的那些宇宙. 但对宇宙而言, 关系 $V\in V$ 不可能成立: 事实上, 那样对 $V$ 的每个子集 $A$ 就都会有 $A\in V$ (\hyperlink{sga4.exp.I.appendix.II.sec.1}{第 1 节}, \hyperlink{sga4.exp.I.appendix.II.sec.1.proposition.1}{命题 1}), 从而 $\card(\mathcal{P}(V))\leq\card(V)$, 这是不可能的.

% END INCLUDED SGA4_English_current_buildable_source_20260906/source/01/drafts/SGA4_I_LINES_7994_8025_English.texfrag


% BEGIN INCLUDED SGA4_English_current_buildable_source_20260906/source/01/drafts/SGA4_I_LINES_8026_8095_English.texfrag
\section{宇宙与强不可达基数}%
\hypertarget{sga4.exp.I.appendix.II.sec.5}{}\pageoriginale

设 $U$ 为宇宙. 由\hyperlink{sga4.exp.I.appendix.II.sec.1.definition.1}{定义 1} 的条件 (U.I), $U$ 的每个元素 $x$ 都是 $U$ 的子集; 因而 $\card(x)\leq\card(U)$. 由于 $U$ 的子集的基数构成一个良序集, 基数
\hypertarget{sga4.exp.I.appendix.II.sec.5.equation.1.cardinal.c.U}{}
\begin{equation}
\sheaf{\sheaf{c}}(U)={\sup}_{x\in U}\card(x)\tag{1}
\end{equation}
存在. 注意对每个基数 $\sheaf{\sheaf{c}}<\sheaf{\sheaf{c}}(U)$, 都存在 $U$ 的元素 $x$ 使得 $\card(x)=\sheaf{\sheaf{c}}$; 事实上, 按定义存在 $y\in U$ 使得 $\sheaf{\sheaf{c}}\leq\card(y)\leq\sheaf{\sheaf{c}}(U)$, 取 $y$ 的适当子集为 $x$ 即可. 反之, 若 $x\in U$, 则 $\card(x)<\sheaf{\sheaf{c}}(U)$; 事实上, 由 \hyperlink{sga4.exp.I.appendix.II.sec.1.definition.1.condition.U.III}{(U.III)} 有 $\mathcal{P}(x)\in U$, 从而 $2^{\card(x)}\leq\sheaf{\sheaf{c}}(U)$. 因此基数 $\sheaf{\sheaf{c}}(U)$ 有下述性质:

\begin{enumerate}
\renewcommand{\labelenumi}{\theenumi)}
\item 若 $\sheaf{\sheaf{c}}$ 为基数且 $\sheaf{\sheaf{c}}<\sheaf{\sheaf{c}}(U)$, 则 $2^{\sheaf{\sheaf{c}}}<\sheaf{\sheaf{c}}(U)$; 事实上, 若 $x$ 为 $U$ 中势为 $\sheaf{\sheaf{c}}$ 的元素, 则由 \hyperlink{sga4.exp.I.appendix.II.sec.1.definition.1.condition.U.III}{(U.III)} 有 $\mathcal{P}(x)\in U$, 从而 $2^{\sheaf{\sheaf{c}}}<\sheaf{\sheaf{c}}(U)$.

\item 若 $(\sheaf{\sheaf{c}}_\lambda)_{\lambda\in I}$ 为基数族, 对每个 $\lambda\in I$ 都有 $\sheaf{\sheaf{c}}_\lambda<\sheaf{\sheaf{c}}(U)$, 且 $\card(I)<\sheaf{\sheaf{c}}(U)$, 则和基数 $\Sigma_{\lambda\in I}\sheaf{\sheaf{c}}_\lambda$ 为 $<\sheaf{\sheaf{c}}(U)$. 事实上, 取 $x_\lambda\in U$ 使得 $\card(x_\lambda)=\sheaf{\sheaf{c}}_\lambda$; 必要时将 $I$ 换成一个等势的指标集, 可假设 $I\in U$; 此时 $x_\lambda$ 的和集是 $U$ 的元素 (\hyperlink{sga4.exp.I.appendix.II.sec.1}{第 1 节}\hyperlink{sga4.exp.I.appendix.II.sec.1.proposition.4}{命题 4} 的\hyperlink{sga4.exp.I.appendix.II.sec.1.corollary.to.proposition.4.2}{推论 2}), 这就证明了断言.
\end{enumerate}

作如下定义:

\hypertarget{sga4.exp.I.appendix.II.sec.5.definition.2}{}
\begin{enonce*}{定义 2}
基数\pageoriginale $\sheaf{\sheaf{d}}$ 称为强不可达的, 如果:
\begin{itemize}
\hypertarget{sga4.exp.I.appendix.II.sec.5.definition.2.condition.FI.1}{}
\item[(FI.1)] 若 $\sheaf{\sheaf{c}}$ 为基数且 $\sheaf{\sheaf{c}}<\sheaf{\sheaf{d}}$, 则 $2^{\sheaf{\sheaf{c}}}<\sheaf{\sheaf{d}}$.

\hypertarget{sga4.exp.I.appendix.II.sec.5.definition.2.condition.FI.2}{}
\item[(FI.2)] 若 $(\sheaf{\sheaf{c}}_\lambda)_{\lambda\in I}$ 为基数族, 对每个 $\lambda\in I$ 都有 $\sheaf{\sheaf{c}}_\lambda<\sheaf{\sheaf{d}}$, 且 $\card(I)<\sheaf{\sheaf{d}}$, 则 $\Sigma_{\lambda\in I}\sheaf{\sheaf{c}}_\lambda<
\sheaf{\sheaf{d}}$.
\end{itemize}
\end{enonce*}

\begin{enonce*}[remark]{例}
基数 $0$ 与可数无限基数都是强不可达的. 非零有限基数都不是强不可达的. 因此我们刚刚证明宇宙公理 \hyperlink{sga4.exp.I.appendix.II.sec.4.axiom.A.6}{(A.6)} 蕴涵以下陈述:
\begin{itemize}
\hypertarget{sga4.exp.I.appendix.II.sec.5.axiom.A-prime.6}{}
\item[${\rm (A'.6)}$] \emph{每个基数都严格小于某个强不可达基数.}
\end{itemize}
反之:
\end{enonce*}

% END INCLUDED SGA4_English_current_buildable_source_20260906/source/01/drafts/SGA4_I_LINES_8026_8095_English.texfrag


% BEGIN INCLUDED SGA4_English_current_buildable_source_20260906/source/01/drafts/SGA4_I_LINES_8096_8244_English.texfrag
\hypertarget{sga4.exp.I.appendix.II.sec.5.theorem.1}{}
\begin{enonce*}{定理 %%
1} 陈述 \hyperlink{sga4.exp.I.appendix.II.sec.5.axiom.A-prime.6}{$(A'.6)$} 蕴涵宇宙公理 (A.6).
\end{enonce*}
事实上, 设 $A$ 为集合. 我们须构造一个以 $A$ 为元素的宇宙 $U$. 递归定义集合序列 $(A_n)_{n\geq0}$ 如下:
\begin{itemize}
\hypertarget{sga4.exp.I.appendix.II.sec.5.theorem.1.proof.equation.2}{}
\item[(2)] $A_0=A$, $A_{n+1}=$ $A_n$ 的各元素的并集 $=$ $A_n$ 的各元素的元素所组成的集合.
\end{itemize}
\noindent 置 $B=\bigcup^\infty_{n=0}A_n$. 根据 \hyperlink{sga4.exp.I.appendix.II.sec.5.axiom.A-prime.6}{$(A'.6)$}, 取强不可达基数 $\sheaf{\sheaf{c}}$ 使得 $\card(B)<\sheaf{\sheaf{c}}$.

存在\emph{良序}集 $I$ 使得 $\card(I)=\sheaf{\sheaf{c}}$. 必要时将 $I$ 换成其中势为 $\sheaf{\sheaf{c}}$ 的最小截段, 可假定\emph{每个与 $I$ 不同的 $I$ 的截段, 其势均为 $<\sheaf{\sheaf{c}}$}. 用\pageoriginale $\mathcal{E}$ 表示 \emph{$I$ 的最小元}. 由关于截段的假设可知 $I$ 没有最大元 (否则将它去掉), 从而 $I$ 的每个元素都有后继; 对 $\alpha\in I$, 用 $s(\alpha)$ 表示 $\alpha$ 的\emph{后继}.

现在用\emph{超限递归}定义集合族 $(B_\alpha)_{\alpha\in I}$ 如下:
\hypertarget{sga4.exp.I.appendix.II.sec.5.theorem.1.proof.equation.3}{}
\begin{equation}
\begin{cases}
& B_{\mathcal{E}}=B\\
& B_{s(\alpha)}=B_{\alpha}\cup \mathcal{P}(B_{\alpha})\\
& B_{\alpha}={\displaystyle{\bigcup_{\beta<\alpha}}}B_{\beta}\mbox{
若 $\alpha$ 没有前驱.}
\end{cases}\tag{3}
\end{equation}
置 $U=\bigcup_{\alpha\in I}B_\alpha$. 我们将证明 $U$ 就是所求的宇宙.

首先, $A\in U$, 因为 $A=A_\circ$ 是 $B=B_{\mathcal{E}}$ 的子集, 因此是 $\mathcal{P}(B_{\mathcal{E}})\subset B_{s(\mathcal{E})}$ 的元素.

其次注意, 对每个 $\alpha\in I$, 有:
\hypertarget{sga4.exp.I.appendix.II.sec.5.theorem.1.proof.equation.4}{}
\begin{equation}
\card(B_{\alpha})<\sheaf{\sheaf{c}}{\quoi}.\tag{4}
\end{equation}
事实上, 对 $\alpha$ 作超限递归. 当 $\alpha=\mathcal{E}$ 时, 按假设成立. 从 \hyperlink{sga4.exp.I.appendix.II.sec.5.theorem.1.proof.equation.4}{(4)} 可得 $\card(B_{s(\alpha)})\leq\card(B_\alpha)+2^{\card(B_\alpha)}<\sheaf{\sheaf{c}}+\sheaf{\sheaf{c}}$ (由 \hyperlink{sga4.exp.I.appendix.II.sec.5.definition.2.condition.FI.1}{(FI.1)}) $=\sheaf{\sheaf{c}}$ (因为 $\sheaf{\sheaf{c}}$ 是无限基数). 最后, 若 $\alpha$ 没有前驱, 则由构造可知 $\beta<\alpha$ 所成的集 $I'$ 的势为 $<\sheaf{\sheaf{c}}$; 因而 (若对每个 $\beta<\alpha$ 有 $\card(B_\beta)<\sheaf{\sheaf{c}}$), $\card(B_\alpha)=\card(\bigcup_{\beta\in I'}B_\beta)\leq\Sigma_{\beta\in I'}\card(B_\beta)<\sheaf{\sheaf{c}}$ (由 (FI.2)). 这就证明了 \hyperlink{sga4.exp.I.appendix.II.sec.5.theorem.1.proof.equation.4}{(4)}.

现在\pageoriginale 证明
\hypertarget{sga4.exp.I.appendix.II.sec.5.theorem.1.proof.equation.5}{}
\begin{equation}
\card(x)<\sheaf{\sheaf{c}}\mbox{ 对每个 }x\in U{\quoi}.\tag{5}
\end{equation}
事实上, 只须证明对每个 $\alpha\in I$ 都有「对每个 $x\in B_\alpha$, $\card(x)<\sheaf{\sheaf{c}}$」. 再对 $\alpha$ 作超限递归. 当 $\alpha=\mathcal{E}$ 时成立, 因为若 $x\in B_{\mathcal{E}}=B$, 则存在 $n\geq0$ 使得 $x\in A_n$; 因而由 \hyperlink{sga4.exp.I.appendix.II.sec.5.theorem.1.proof.equation.2}{(2)} 有 $x\subset A_{n+1}$, $x\subset B$, 且 $\card(x)\leq\card(B)<\sheaf{\sheaf{c}}$. $\alpha$ 没有前驱的情形显然. 最后, 若 $x\in B_\alpha$ 且 $\alpha=s(\beta)$, 则或者 $x\in B_\beta$, 此时由递归假设断言「$\card(x)<\sheaf{\sheaf{c}}$」成立; 或者 $x\in\mathcal{P}(B_\beta)$, 从而 $x\subset B_\beta$, 此时由 \hyperlink{sga4.exp.I.appendix.II.sec.5.theorem.1.proof.equation.4}{(4)} 知断言「$\card(x)<\sheaf{\sheaf{c}}$」成立.

现在可以证明 $U$ 确是一个\emph{宇宙}:

\begin{itemize}
\hypertarget{sga4.exp.I.appendix.II.sec.5.theorem.1.proof.clause.U.I}{}
\item[(U.I)] 设 $x\in U$ 且 $y\in x$. 须证明 $y\in U$. 换言之, 须证明对每个 $\alpha\in I$, 有:
$$
\textrm{``}x\in B_{\alpha}\mbox{ 且 }y\in x\Longrightarrow y\in
B_{\alpha}\textrm{''\,}{\quoi}.
$$
再次对 $\alpha$ 作超限递归. 当 $\alpha=\mathcal{E}$ 时成立, 因为 $x\in B_{\mathcal{E}}=B$ 蕴涵对某个 $n$ 有 $x\in A_n$; 因而若 $y\in x$, 则 $y\in A_{n+1}$, 从而 $y\in B=B_{\mathcal{E}}$. 推至没有前驱的元素 $\alpha$ 是显然的. 最后, 从 $\alpha$ 推至 $s(\alpha)$: 若 $x\in B_{s(\alpha)}$ 且 $y\in x$, 则或者 $x\in B_\alpha$, 从而由递归假设有 $y\in B_\alpha\subset B_{s(\alpha)}$; 或者 $x\in\mathcal{P}(B_\alpha)$, 从而仍有 $y\in B_\alpha\subset B_{s(\alpha)}$.

\hypertarget{sga4.exp.I.appendix.II.sec.5.theorem.1.proof.clause.U.II}{}
\item[(U.II)] 设 $x,y\in U$. 由 \hyperlink{sga4.exp.I.appendix.II.sec.5.theorem.1.proof.equation.3}{(3)} 可知族 $(B_\alpha)_{\alpha\in I}$ 递增, 因而存在 $\alpha\in I$ 使得 $x,y\in B_\alpha$. 此时 $\{x,y\}\in\mathcal{P}(B_\alpha)\subset B_{s(\alpha)}\subset U$.

\hypertarget{sga4.exp.I.appendix.II.sec.5.theorem.1.proof.clause.U.III}{}
\item[(U.III)] 设\pageoriginale $x\in U$. 存在 $\alpha\in I$ 使得 $x\in B_\alpha$. 因此只须证明, 对每个 $\alpha\in I$, 有:
$$
\textrm{``}x\in B_{\alpha}\Longrightarrow \mathcal{P}(x)\in
B_{s(s(\alpha))}\textrm{''\,}{\quoi}.
$$
再次对 $\alpha$ 作超限递归. 若 $\alpha=\mathcal{E}$, 则对某个 $n$ 有 $x\in A_n$, 从而对 $x$ 的每个子集 $y$, 有 $y\subset A_{n+1}\subset B=B_{\mathcal{E}}$; 因而对每个 $y\in\mathcal{P}(x)$, 有 $y\in\mathcal{P}(B_{\mathcal{E}})\subset B_{s(\mathcal{E})}$, 从而 $\mathcal{P}(x)\subset B_{s(\mathcal{E})}$, 且 $\mathcal{P}(x)\in\mathcal{P}(B_{s(\mathcal{E})})\subset B_{s(s(\mathcal{E}))}$, 故当 $\alpha=\mathcal{E}$ 时断言成立. 推至没有前驱的元素 $\alpha$ 是直接的, 因为此时 $\beta<\alpha$ 蕴涵 $s(s(\beta))<\alpha$. 最后, 从 $\alpha$ 推至 $s(\alpha)$; 设 $x\in B_{s(\alpha)}$; 若 $x\in B_\alpha$, 则由递归假设有 $\mathcal{P}(x)\in B_{s(s(\alpha))}\subset B_{s(s(s(\alpha)))}$; 若 $x\in\mathcal{P}(B_\alpha)$, 则 $x\subset B_\alpha$, 从而 $\mathcal{P}(x)\subset\mathcal{P}(B_\alpha)$, 且 $\mathcal{P}(x)\in\mathcal{P}(\mathcal{P}(B_\alpha))\in B_{s(s(s(\alpha)))}$.

\hypertarget{sga4.exp.I.appendix.II.sec.5.theorem.1.proof.clause.U.IV}{}
\item[(U.IV)] 设 $(x_\lambda)_{\lambda\in K}$ 为 $U$ 的元素族, 且 $K\in U$. 须证明并集 $x=\bigcup_{\lambda\in K}x_\lambda$ 是 $U$ 的元素. 对每个 $\lambda\in K$, 取 $\alpha(\lambda)\in I$ 使得 $x_\lambda\in B_{\alpha(\lambda)}$. 来证明各 $\alpha(\lambda)$ 所成的集合 $\alpha(K)$ 在 $I$ 中\emph{有上界}: 事实上, 否则就有
$$
I=\bigcup_{\lambda%%
\in K}[\mathcal{E},\alpha(\lambda)]{\quoi},
$$
这将与 \hyperlink{sga4.exp.I.appendix.II.sec.5.definition.2.condition.FI.2}{(FI.2)} 矛盾, 因为 $\card(K)<\sheaf{\sheaf{c}}$, 且对每个 $\lambda\in K$ 有 $\card([\mathcal{E},\alpha(\lambda)])<\sheaf{\sheaf{c}}$. 于是设 $\beta$ 为 $\alpha(K)$ 的一个上界; 因为族 $(B_\alpha)_{\alpha\in I}$ 递增, 对每个 $\lambda\in K$ 都有 $x_\lambda\in B_\beta$. 由 \hyperlink{sga4.exp.I.appendix.II.sec.5.theorem.1.proof.clause.U.I}{(U.I)} 中所证, $x_\lambda$ 的每个元素都属于 $B_\beta$, 从而 $x=\bigcup_{\lambda\in K}x_\lambda\subset B_\beta$. 由 \hyperlink{sga4.exp.I.appendix.II.sec.5.theorem.1.proof.equation.3}{(3)}, 得 $x\in B_{s(\beta)}$, 因而 $x\in U$. \cqfd
\end{itemize}

% END INCLUDED SGA4_English_current_buildable_source_20260906/source/01/drafts/SGA4_I_LINES_8096_8244_English.texfrag


% BEGIN INCLUDED SGA4_English_current_buildable_source_20260906/source/01/drafts/SGA4_I_LINES_8245_8358_English.texfrag
\hypertarget{sga4.exp.I.appendix.II.sec.5.definition.3}{}
\begin{enonce*}{定义 3}
设\pageoriginale $x$, $y$ 为两个集合, $n$ 为整数 $\geq 0$. 如果存在序列 $(x_j)_{j=0,\ldots,n}$, 使得 $x_0=x$, $x_n=y$, 且对 $j=0,\ldots,n-1$ 有 $x_{j+1}\in x_j$, 就说 $y$ 是 $x$ 的 $n$ 阶分量.
\end{enonce*}

因而 $x$ 是 $x$ 唯一的 0 阶分量. $x$ 的 1 阶 (分别为 2 阶) 分量就是 $x$ 的元素 (分别为 $x$ 的元素的元素). 如果存在 $n\geq 0$ 使得 $y$ 是 $x$ 的 $n$ 阶分量, 就说 $y$ 是 $x$ 的\emph{分量}. 关系「$y$ 是 $x$ 的分量」是一个\emph{预序关系}. 由选择--并集公理模式 (第 II 章, \ldots), 关系「$y$ 是 $x$ 的分量」关于 $y$ 可集化, 因而 $x$ 的分量构成一个\emph{集合}.

\hypertarget{sga4.exp.I.appendix.II.sec.5.definition.4}{}
\begin{enonce*}{定义 4}
设 $\sheaf{\sheaf{c}}$ 为基数. 如果集合 $x$ 的所有分量的势都为 $\leq \sheaf{\sheaf{c}}$ (分别为 $<\sheaf{\sheaf{c}}$、有限势), 就说 $x$ 是 $\sheaf{\sheaf{c}}$ 型的 (分别为严格 $\sheaf{\sheaf{c}}$ 型的、有限型的).
\end{enonce*}

\begin{enonce*}[remark]{例}
宇宙 $U_1$ 的元素 (\hyperlink{sga4.exp.I.appendix.II.sec.1}{第 1 节}, \hyperlink{sga4.exp.I.appendix.II.sec.1.example.2}{例 2}) 都是有限型的. 若 $\sheaf{\sheaf{c}}$ 为基数, 且 $x$ 是 $\sheaf{\sheaf{c}}$ 型的 (分别为严格 $\sheaf{\sheaf{c}}$ 型的、有限型的), 则 $x$ 的每个分量及每个子集都是 $\sheaf{\sheaf{c}}$ 型的 (分别为严格 $\sheaf{\sheaf{c}}$ 型的、有限型的); 同样, $\mathcal{P}(x)$ 是 $2^{\sheaf{\sheaf{c}}}$ 型的 (分别为严格 $2^{\sheaf{\sheaf{c}}}$ 型的、有限型的). 若 $x$ 为 $\sheaf{\sheaf{c}}$ 型且 $\sheaf{\sheaf{c}}\leq\sheaf{\sheaf{c}}'$, 则 $x$ 为 $\sheaf{\sheaf{c}}'$ 型.
\end{enonce*}

\hypertarget{sga4.exp.I.appendix.II.sec.5.lemma.of.section.5}{}
\begin{enonce*}{引理}
设\pageoriginale $\sheaf{\sheaf{c}}$ 为不可数强不可达基数, $X$ 为严格 $\sheaf{\sheaf{c}}$ 型的集合. 则存在基数 $\sheaf{\sheaf{d}}<\sheaf{\sheaf{c}}$, 使得 $X$ 为 $\sheaf{\sheaf{d}}$ 型.
\end{enonce*}

事实上, 对每个整数 $n$, 用 $\sheaf{\sheaf{d}}_n$ 表示 $X$ 的 $n$ 阶分量集 $X_n$ 的势. 有 $\sheaf{\sheaf{d}}_0=1$, $\sheaf{\sheaf{d}}_1=\card(X)<\sheaf{\sheaf{c}}$, 且由于 $X_n=\bigcup_{y\in X_{n-1}} y$, 对 $n$ 作归纳并用 \hyperlink{sga4.exp.I.appendix.II.sec.5.definition.2.condition.FI.2}{(FI.2)}, 得对每个 $n$ 有 $\sheaf{\sheaf{d}}_n<\sheaf{\sheaf{c}}$. 因而各 $\sheaf{\sheaf{d}}_n$ 以基数 $\sheaf{\sheaf{d}}=\Sigma_{j\geq 0}\sheaf{\sheaf{d}}_j$ 为上界, 由 (FI.2) 和 $\sheaf{\sheaf{c}}$ 不可数的假设, 此基数为 $<\sheaf{\sheaf{c}}$. 若 $Y$ 为 $X$ 的 $n$ 阶分量, 则 $\card(Y)\leq \card(X_{n+1})\leq\sheaf{\sheaf{d}}$. \cqfd

\hypertarget{sga4.exp.I.appendix.II.sec.5.proposition.10}{}
\begin{enonce*}{命题 10}
设 $U$ 为宇宙, $\sheaf{\sheaf{c}}$ 为强不可达基数. 则由 $U$ 中严格 $\sheaf{\sheaf{c}}$ 型的那些 $x\in U$ 构成的集合 $U'$ 是宇宙.
\end{enonce*}
事实上, 若 $x$, $y\in U'$ 且 $z\in x$, 则显然 $z\in U'$ 且 $\{x,y\}\in U'$. 若 $x\in U'$, 则 $\card(x)<\sheaf{\sheaf{c}}$, 从而由 \hyperlink{sga4.exp.I.appendix.II.sec.5.definition.2.condition.FI.1}{(FI.1)} 得 $\card(\mathcal{P}(x))<\sheaf{\sheaf{c}}$; 因为 $\mathcal{P}(x)$ 的分量或者是 $\mathcal{P}(x)$ 本身, 或者是 $x$ 的子集, 或者是 $x$ 的分量, 故有 $\mathcal{P}(x)\in U'$. 最后, 若 $(x_{\alpha})_{\alpha\in I}$ 为 $U'$ 的元素族, 且 $I\in U'$, 则由 \hyperlink{sga4.exp.I.appendix.II.sec.5.definition.2.condition.FI.2}{(FI.2)} 得 $\card(\bigcup_{\alpha}x_{\alpha})<\sheaf{\sheaf{c}}$; 因为 $\bigcup_{\alpha}x_{\alpha}$ 的每个 $>0$ 阶分量都是某个 $x_{\alpha}$ 的分量, 确有 $\bigcup_{\alpha}x_{\alpha}\in U'$. \cqfd

\medskip
\noindent \emph{关于基数} $\sheaf{\sheaf{c}}(U)$ \emph{的注}.
设 $U$ 为宇宙, $\sheaf{\sheaf{c}}(U)$ 为 \hyperlink{sga4.exp.I.appendix.II.sec.5.equation.1.cardinal.c.U}{(1)} 定义的基数, 即
$$
\sheaf{\sheaf{c}}(U)={\sup}_{x\in U}\card(x){\quoi}.
$$
显然\pageoriginale $\sheaf{\sheaf{c}}(U)\leq \card(U)$, 因为回想 $x\in U$ 蕴涵 $x\subset U$. 但等式 $\sheaf{\sheaf{c}}(U)=\card(U)$ 并非总成立. 例如, 设 $(x_{\alpha})_{\alpha\in I}$ 为满足下述公理的\emph{不可数}符号族:
$$
\text{«} x_{\alpha}=\{x_{\alpha}\}\mbox{ 对每个 } \alpha\in I
\text{»} \quav \text{«} x_{\alpha}\neq x_{\beta}\mbox{ 若 }
\alpha\neq \beta \text{»}
$$
(换言之, $x_{\alpha}$ 是以自身为唯一元素的单元素集). 如同\hyperlink{sga4.exp.I.appendix.II.sec.1}{第 1 节}\hyperlink{sga4.exp.I.appendix.II.sec.1.example.2}{例 2}, 用符号 $∅$, $x_{\alpha}$, $\{,\}$ 等构造「有意义的字」. 每个字都表示\emph{有限}集. 设 $U$ 为这些字的集合. 容易验证\hyperlink{sga4.exp.I.appendix.II.sec.1.definition.1}{定义 1} 的条件成立, 因而 $U$ 为\emph{宇宙}. 因为有意义的字是势为 $\card(I)+4=\card(I)$ 的集合的元素的有限序列, 所以 $\card(U)=\card(I)$ (第 III 章; 其实 $\card(U)\geq \card(I)$ 即已足够, 而这显然成立). 另一方面, $\sheaf{\sheaf{c}}(U)$ 为可数无限基数, 因而 $\sheaf{\sheaf{c}}(U)<\card(U)$.

严格不等式 $\sheaf{\sheaf{c}}(U)<\card(U)$ 源于我们引入了满足 $x_{\alpha}\in x_{\alpha}$ 的集合 $x_{\alpha}$. 如果排除这类怪物, 就得到对每个宇宙 $U$ 都有 $\card(U)=\sheaf{\sheaf{c}}(U)$, 以及其他一些十分漂亮的「不可能有什么用处」的结果. 下一节我们就来做这件事.

% END INCLUDED SGA4_English_current_buildable_source_20260906/source/01/drafts/SGA4_I_LINES_8245_8358_English.texfrag


% BEGIN INCLUDED SGA4_English_current_buildable_source_20260906/source/01/drafts/SGA4_I_LINES_8359_8513_English.texfrag
\section{Artin 集与宇宙}%
\hypertarget{sga4.exp.I.appendix.II.sec.6}{}

\hypertarget{sga4.exp.I.appendix.II.sec.6.definition.5}{}
\begin{enonce*}{定义 5}
集合 $x$ 称为 Artin 集, 如果不存在无限序列 $(x_n)_{n\geq 0}$, 使得 $x_{\circ}=x$, 且对每个 $n\geq 0$ 有 $x_{n+1}\in x_n$.
\end{enonce*}

\begin{enonce*}[remark]{例}
集合 $∅$, $\{∅\}$, $\{∅,\{∅\}\}$, 更一般地, \hyperlink{sga4.exp.I.appendix.II.sec.1}{第 1 节} 的宇宙 $U_1$ 的元素, 都是 Artin 集.
\end{enonce*}

形象地说\pageoriginale, 若 $x$ 为 Artin 集, 在 $x$ 中取元素 $x_1$, 再在 $x_1$ 中取元素 $x_2$, 等等, 这个过程必定停止, 并到达 $x$ 的一个\emph{空的}分量 $x_n$. 换言之, Artin 集是「从 $∅$ 构成的」; 以后将把这一点说得更精确 (参见\hyperlink{sga4.exp.I.appendix.II.sec.6.theorem.2}{定理 2} 证明中的 \hyperlink{sga4.exp.I.appendix.II.sec.6.theorem.2.proof.claim.star}{(*)}).

Artin 集显然有以下性质:
\begin{itemize}
\hypertarget{sga4.exp.I.appendix.II.sec.6.property.AR.I}{}
\item[(AR.I)] 若 $x$ 为 Artin 集, 则 $x$ 的每个子集和每个分量都是 Artin 集. $x$ 为 Artin 集的充分必要条件是 $x$ 的每个元素都是 Artin 集.

\hypertarget{sga4.exp.I.appendix.II.sec.6.property.AR.II}{}
\item[(AR.II)] 若 $x$ 与 $y$ 为 Artin 集, 则 $\{x,y\}$ 为 Artin 集.

\hypertarget{sga4.exp.I.appendix.II.sec.6.property.AR.III}{}
\item[(AR.III)] 若 $x$ 为 Artin 集, 则 $\mathcal{P}(x)$ 为 Artin 集.

\hypertarget{sga4.exp.I.appendix.II.sec.6.property.AR.IV}{}
\item[(AR.IV)] Artin 集的任意并仍是 Artin 集.
\end{itemize}

这些性质立即给出以下结果:

\hypertarget{sga4.exp.I.appendix.II.sec.6.proposition.11}{}
\begin{enonce*}{命题 11}
若 $U$ 为宇宙, 则 $U$ 中 Artin 元素的集合为宇宙, 且必为 Artin 集.
\end{enonce*}

\hypertarget{sga4.exp.I.appendix.II.sec.6.corollary.to.proposition.11}{}
\begin{enonce*}{推论}
若 $x$ 为 Artin 集, 则存在 Artin 宇宙 $U$ 使得 $x\in U$.
\end{enonce*}

事实上, 由公理 \hyperlink{sga4.exp.I.appendix.II.sec.4.axiom.A.6}{(A.6)}, $x$ 是某个宇宙 $V$ 的元素; 取 $V$ 中 Artin 元素的集合为 $U$ 即可.

以下命题甚至比第 6 节的其他内容还要没用:

\hypertarget{sga4.exp.I.appendix.II.sec.6.proposition.12}{}
\begin{enonce*}{命题 12}
设\pageoriginale $A$ 为 Artin 集. 则:
\begin{enumerate}
\renewcommand{\theenumi}{\alph{enumi}}
\renewcommand{\labelenumi}{\theenumi)}
\hypertarget{sga4.exp.I.appendix.II.sec.6.proposition.12.part.a}{}
\item 对每个 $x\in A$, 有 $x\not\in x$;

\hypertarget{sga4.exp.I.appendix.II.sec.6.proposition.12.part.b}{}
\item 若 $x$, $y\in A$, 则 $x\in y$ 和 $y\in x$ 不可能同时成立;

\hypertarget{sga4.exp.I.appendix.II.sec.6.proposition.12.part.c}{}
\item 对每个 $x\in A$, 关系「$y$ 是 $z$ 的分量」是 $x$ 的分量 $y$, $z$ 上的序关系;

\hypertarget{sga4.exp.I.appendix.II.sec.6.proposition.12.part.d}{}
\item 对 $A$ 的每个非空元素 $x$, 存在 $y\in x$ 使得 $x\cap y=∅$.
\end{enumerate}
\end{enonce*}
事实上, a) (分别为 b)) 的否定会给出与\hyperlink{sga4.exp.I.appendix.II.sec.6.definition.5}{定义 5} 矛盾的无限序列 $(x_n)_{n\geq 0}$, 即 $(x,x,x,\ldots)$ (分别为 $(x,y,x,y,x,y,\ldots)$). c) 的否定表示存在 $x\in A$, $x$ 的互异分量 $y$, $z$, 以及隶属关系链:
$$
y\in y_1\in\ldots \in y_q\in z\quav z\in z_1\in\ldots\in z_r\in
y{\quoi};
$$
从而如同 \hyperlink{sga4.exp.I.appendix.II.sec.6.proposition.12.part.b}{b)}, 得到与\hyperlink{sga4.exp.I.appendix.II.sec.6.definition.5}{定义 5} 矛盾的无限序列 $(x_n)_{n\geq 0}$. 最后, 若 d) 为假, 则存在 $A$ 的非空元素 $x$, 使对每个 $y\in x$ 有 $x\cap y\neq ∅$; 置 $x_\circ=x$, 并取 $x$ 的一个元素为 $x_1$; 因为 $x\cap x_1\neq ∅$, 取 $x\cap x_1$ 的一个元素为 $x_2$, 如此继续; 更形式地, 用 $x_1\in x$, 以及对 $n\geq 1$ 用 $x_{n+1}\in x_n\cap x$, 归纳定义由 $x$ 的元素组成的无限序列 $(x_n)_{n\geq 1}$; 此时序列 $(x_n)_{n\geq 0}$ 与\hyperlink{sga4.exp.I.appendix.II.sec.6.definition.5}{定义 5} 矛盾.

\begin{enonce*}[remark]{注}
设 $B$ 为集合. 则 $B$ 为 Artin 集, 当且仅当每个以 $B$ 的分量集为元素的集合 $A$, 都满足\hyperlink{sga4.exp.I.appendix.II.sec.6.proposition.12}{命题 12} 的\hyperlink{sga4.exp.I.appendix.II.sec.6.proposition.12.part.d}{条件 d)}. 事实上, 必要性由 \hyperlink{sga4.exp.I.appendix.II.sec.6.property.AR.I}{(AR.I)} 和\pageoriginale \hyperlink{sga4.exp.I.appendix.II.sec.6.proposition.12}{命题 12} 得出. 反之, 若 $B$ 不是 Artin 集, 则存在无限序列 $(x_n)_{n\geq0}$, 使得 $x_{\circ}=B$, 且对每个 $n\geq 0$ 有 $x_{n+1}\in x_n$; 此时, 取以各 $x_n$ 所成的集合 $X$ 为唯一元素的集合为 $A$; 则 $A$ 含有非空元素 $X$, 且对 $X$ 的每个元素 $y$ 都有 $y\cap X\neq ∅$ (事实上, $y$ 形如 $x_n$, 而 $x_{n+1}\in y\cap X$).
\end{enonce*}

\hypertarget{sga4.exp.I.appendix.II.sec.6.theorem.2}{}
\begin{enonce*}{定理 2}
设 $\sheaf{\sheaf{c}}$ 为无限基数. 则:
\begin{itemize}
\hypertarget{sga4.exp.I.appendix.II.sec.6.theorem.2.part.a}{}
\item[a)] 关系「$x$ 是 $\sheaf{\sheaf{c}}$ 型 (分别为严格 $\sheaf{\sheaf{c}}$ 型) 的 Artin 集」关于 $x$ 可集化; $\sheaf{\sheaf{c}}$ 型的 Artin 集所成的集合 $A_{\sheaf{\sheaf{c}}}$ 的势为 $2^{\sheaf{\sheaf{c}}}$.

\hypertarget{sga4.exp.I.appendix.II.sec.6.theorem.2.part.b}{}
\item[b)] 若 $\sheaf{\sheaf{c}}$ 为强不可达基数, 则严格 $\sheaf{\sheaf{c}}$ 型的 Artin 集所成的集合 $U_{\sheaf{\sheaf{c}}}$ 是势为 $\sheaf{\sheaf{c}}$ 的宇宙; 基数 $\sheaf{\sheaf{c}}(U_{\sheaf{\sheaf{c}}})=\sup_{x\in U_{\sheaf{\sheaf{c}}}}\card(x)$ 为 $\sheaf{\sheaf{c}}$.

\hypertarget{sga4.exp.I.appendix.II.sec.6.theorem.2.part.c}{}
\item[c)] 若宇宙 $U$ 含有势为 $\sheaf{\sheaf{c}}$ 的元素, 则每个 $\sheaf{\sheaf{c}}$ 型的 Artin 集都属于 $U$ (换言之, $A_{\sheaf{\sheaf{c}}}\subset U$).

\hypertarget{sga4.exp.I.appendix.II.sec.6.theorem.2.part.c-prime}{}
\item[${\rm c')}$] 若宇宙 $U$ 非空, 则每个有限型的 Artin 集都是 $U$ 的元素.
\end{itemize}
\end{enonce*}
在证明\hyperlink{sga4.exp.I.appendix.II.sec.6.theorem.2}{定理 2} 之前, 先由它推出几个有启发性的推论:

\hypertarget{sga4.exp.I.appendix.II.sec.6.corollary.to.theorem.2.1}{}
\begin{enonce*}{推论 1}
若宇宙 $U$ 为 Artin 集, 则 $\card(U)$ 为强不可达基数, 且 $U$ 是严格 $\card(U)$ 型的 Artin 集所成的集合.
\end{enonce*}

事实上,\pageoriginale 置 $\sheaf{\sheaf{c}}(U)=\sup_{x\in U}\card(x)$. 这是强不可达基数 (\hyperlink{sga4.exp.I.appendix.II.sec.5}{第 5 节}开头). 先假定它不可数; 则对每个无限基数 $\sheaf{\sheaf{c}}<\sheaf{\sheaf{c}}(U)$, 由 \hyperlink{sga4.exp.I.appendix.II.sec.6.theorem.2.part.c}{c)} 可知每个 $\sheaf{\sheaf{c}}$ 型的 Artin 集都是 $U$ 的元素; 因而由\hyperlink{sga4.exp.I.appendix.II.sec.5.lemma.of.section.5}{第 5 节的引理}, 每个严格 $\sheaf{\sheaf{c}}(U)$ 型的 Artin 集都是 $U$ 的元素. 若 $\sheaf{\sheaf{c}}(U)$ 可数, 则由 \hyperlink{sga4.exp.I.appendix.II.sec.6.theorem.2.part.c-prime}{${\rm c'}$)} 可知这个断言仍然成立. 反之, 由 (U.I), $U$ 中的每个集合都是严格 $\sheaf{\sheaf{c}}(U)$ 型. 因此 $U$ 是 \hyperlink{sga4.exp.I.appendix.II.sec.6.theorem.2.part.b}{b)} 中的宇宙 $U_{\sheaf{\sheaf{c}}(U)}$, 从而由 b) 有 $\sheaf{\sheaf{c}}(U)=\card(U)$.

由\hyperlink{sga4.exp.I.appendix.II.sec.6.corollary.to.theorem.2.1}{推论 1} 可知, 一个 Artin 宇宙由它的势 (此势并且是强不可达基数) 唯一\emph{确定}. 因而 Artin 宇宙与强不可达基数之间存在「一一对应」. 特别地:

% END INCLUDED SGA4_English_current_buildable_source_20260906/source/01/drafts/SGA4_I_LINES_8359_8513_English.texfrag


% BEGIN INCLUDED SGA4_English_current_buildable_source_20260906/source/01/drafts/SGA4_I_LINES_8514_8642_English.texfrag
\hypertarget{sga4.exp.I.appendix.II.sec.6.corollary.to.theorem.2.2}{}
\begin{enonce*}{推论 2}
Artin 宇宙上的包含关系 $U\subset U'$ 是良序关系.
\end{enonce*}

事实上, 基数上的关系 $\sheaf{\sheaf{c}}\leq \sheaf{\sheaf{c}}'$ 是良序关系 (第 III 章), 而用\hyperlink{sga4.exp.I.appendix.II.sec.6.theorem.2}{定理 2} 的 \hyperlink{sga4.exp.I.appendix.II.sec.6.theorem.2.part.b}{b)} 中的记号, 关系 $\sheaf{\sheaf{c}}\leq\sheaf{\sheaf{c}}'$ 与 $U_{\sheaf{\sheaf{c}}}\subset U_{\sheaf{\sheaf{c}}'}$ 等价.

注意\hyperlink{sga4.exp.I.appendix.II.sec.6.theorem.2}{定理 2} 的 \hyperlink{sga4.exp.I.appendix.II.sec.6.theorem.2.part.b}{b)} 给出了\hyperlink{sga4.exp.I.appendix.II.sec.5.theorem.1}{定理 1} (\hyperlink{sga4.exp.I.appendix.II.sec.5}{第 5 节}) 的另一证明.

现在来证明\hyperlink{sga4.exp.I.appendix.II.sec.6.theorem.2}{定理 2}. 给定集合 $A$, 称满足 $x_0=A$, 且对 $j=0,\ldots,n-1$ 有 $x_{j+1}\in x_j$ 的任意有限序列 $(x_j)_{j=0,\ldots,n}$ 为 $A$ 的\emph{链}. 由选择--并集公理模式, $A$ 的链构成一个\pageoriginale \emph{集合}, 记为 $G(A)$. 给定链 $X=(x_j)_{j=0,\ldots,n}$, 称形如 $(x_i)_{i=0,\ldots,q}$ 且 $q\leq n$ 的链小于 $X$; 这在 $G(A)$ 上赋予了\emph{有序集}结构. $A$ 为 Artin 集的充分必要条件是 $G(A)$ 为「\emph{Noether}」有序集 (即满足第 III 章, \S 6, 第 5 节中的等价条件).

我们将证明:

\hypertarget{sga4.exp.I.appendix.II.sec.6.theorem.2.proof.claim.star}{}
\noindent (*) \emph{若 $A$ 为 Artin 集, 则它由有序集 $G(A)$ 的同构类唯一确定.}

事实上, 给定有序集 $G$ 及元素 $g\in G$, 用 $S(g)$ 表示满足 $g<g'$, 且 $g\leq h\leq g'$ 蕴涵 $h=g$ 或 $h=g'$ 的一切 $g'\in G$ 所成的集合 (换言之, $g$ 的「直接后继」集). 考虑将 $G(A)$ 中的每个链 $X=(x_0,\ldots,x_n)$ 映为集合 $x_n$ 的映射 $\theta$; 则对每个 $X\in G(A)$, 有:
$$
\theta(X)=\{\theta(X')| X'\in S(X)\}{\quoi}.
$$
因为 $A$ 是 $G(A)$ 的最小元 $X_0=(A)$ 在 $\theta$ 下的像, 只须证明 $\theta$ 由有序集 $G(A)$ 的同构类唯一确定. 这由以下引理得出:

\hypertarget{sga4.exp.I.appendix.II.sec.6.lemma.1}{}
\begin{enonce*}{引理 1}
设\pageoriginale $G$ 为 Noether 有序集, $\varphi$ 为定义域是 $G$ 的函数, 使对每个 $g\in G$, 有:
\hypertarget{sga4.exp.I.appendix.II.sec.6.lemma.1.equation.1}{}
\begin{equation}
\varphi(g)=\{\varphi(g')| g'\in S(g)\}{\quoi}.\tag{1}
\end{equation}
则 $\varphi$ 唯一确定. 此外, 若 $U$ 为包含与 $G$ 等势的元素的宇宙, 则 $\varphi$ 在 $U$ 中取值.
\end{enonce*}

假设蕴涵: 若 $g$ 是 $G$ 的极大元, 则 $\varphi(g)=∅$.

事实上, 设 $\varphi$ 与 $\varphi'$ 为满足 \hyperlink{sga4.exp.I.appendix.II.sec.6.lemma.1.equation.1}{(1)} 的两个函数. 若 $\varphi\neq \varphi'$, 则满足 $\varphi(g)\neq \varphi'(g)$ 的一切 $g\in G$ 所成的集合非空, 因而因 $G$ 为 Noether 集, 它有一个\emph{极大}元 $h$. 此时对每个 $g>h$ 有 $\varphi(g)=\varphi'(g)$, 特别地对每个直接后继 $g\in S(h)$ 成立; 因而由 \hyperlink{sga4.exp.I.appendix.II.sec.6.lemma.1.equation.1}{(1)} 有 $\varphi(h)=\varphi'(h)$, 矛盾. 故确有 $\varphi=\varphi'$. 现在设 $U$ 为包含与 $G$ 等势的元素 $x$ 的宇宙; 来证明 $\varphi$ 在 $U$ 中取值. 否则, 在满足 $\varphi(g)\not\in U$ 的各 $g\in G$ 中取极大元 $h$. 则对每个 $g'\in S(h)$, 有 $\varphi(g')\in U$, 因而
$$
\varphi(h)=\{\varphi(g')| g'\in S(h)\}
$$
是 $U$ 的子集. 因为它的势不超过 $\card(G)$, 从而不超过 $U$ 的某个元素的势, 故 $\varphi(h)\in U$ (\hyperlink{sga4.exp.I.appendix.II.sec.1}{第 1 节}, \hyperlink{sga4.exp.I.appendix.II.sec.1.proposition.7}{命题 7}). 这个矛盾表明 $\varphi$ 在 $U$ 中取值.

由此\pageoriginale 可证明\hyperlink{sga4.exp.I.appendix.II.sec.6.theorem.2}{定理 2} 的 a). 只须证明不带「分别」部分的断言, 因为另一断言随即得出. 设 $\sheaf{\sheaf{c}}$ 为\emph{无限}基数. 若 $A$ 为 $\sheaf{\sheaf{c}}$ 型的集合, 则:
\hypertarget{sga4.exp.I.appendix.II.sec.6.theorem.2.proof.equation.2}{}
\begin{equation}
\card(G(A))\leq \sheaf{\sheaf{c}}{\quoi}%%
\ .\tag{2}
\end{equation}
事实上, 若 $A_n$ 表示 $A$ 的 $n$ 阶分量所成的集合, 则 $\card(A_1)\leq \sheaf{\sheaf{c}}$, 且 $\card(A_{n+1})\leq\sheaf{\sheaf{c}}\cdot \card(A_n)$, 从而归纳得 $\card(A_n)\leq\sheaf{\sheaf{c}}^n$. 但因为 $\sheaf{\sheaf{c}}$ 无限, 有 $\sheaf{\sheaf{c}}^n=\sheaf{\sheaf{c}}$ (第 III 章); 因而 $\card(\bigcup^\infty_{n=0} A_n)\leq\card(\NN)\cdot \sheaf{\sheaf{c}}=\sheaf{\sheaf{c}}$ (第 III 章). 而 $G(A)$ 是 $\bigcup_n A_n$ 的元素的\emph{有限}序列所成的集合, 因此不等式 (2) 确实成立 (第 III 章). 现在, 若 $E$ 为势为 $\sheaf{\sheaf{c}}$ 的集合, 在 $E$ 的一个子集 $E'$ 上给定序结构, 等价于给定由满足 $x\in E'$, $y\in E'$ 且 $x\leq y$ 的 $(x,y)$ 组成的 $E\times E$ 的子集. 因此, 由 \hyperlink{sga4.exp.I.appendix.II.sec.6.theorem.2.proof.equation.2}{(2)}, 有序集 $G(A)$ (其中 $A$ 为 $\sheaf{\sheaf{c}}$ 型) 的同构类构成集合 $\mathfrak{F}_{\sheaf{\sheaf{c}}}$, 且 $\card(\mathfrak{F}_{\sheaf{\sheaf{c}}})\leq\card(\mathcal{P}(E\times E))=2^{\sheaf{\sheaf{c}}}$. 设 $\mathfrak{F}'_{\sheaf{\sheaf{c}}}$ 为 $\mathfrak{F}_{\sheaf{\sheaf{c}}}$ 中由具有最小元的 Noether 有序集的类组成的子集. 若对每个 $G\in\mathfrak{F}'_{\sheaf{\sheaf{c}}}$, 配以\hyperlink{sga4.exp.I.appendix.II.sec.6.lemma.1}{引理 1} 的函数 $\varphi$ 在 $G$ 的最小元处的值, 就得到 $\mathfrak{F}'_{\sheaf{\sheaf{c}}}$ 上的映射 $\theta$, 它的像包含\emph{一切} $\sheaf{\sheaf{c}}$ 型的 Artin 集. 因而后者确实构成集合 $A_{\sheaf{\sheaf{c}}}$, 且有 $\card(A_{\sheaf{\sheaf{c}}})\leq\card(\mathfrak{F}'_{\sheaf{\sheaf{c}}})\leq\card(\mathfrak{F}_{\sheaf{\sheaf{c}}})\leq2^{\sheaf{\sheaf{c}}}$.

还须\pageoriginale 证明 $\card(A_{\sheaf{\sheaf{c}}})=2^{\sheaf{\sheaf{c}}}$. 为此, 只须证明存在一个 $\sheaf{\sheaf{c}}$ 型且势为 $\sheaf{\sheaf{c}}$ 的 Artin 集 $B$, 因为此时 $B$ 的子集都会是 $A_{\sheaf{\sheaf{c}}}$ 的元素. 此存在性由以下引理得出:

% END INCLUDED SGA4_English_current_buildable_source_20260906/source/01/drafts/SGA4_I_LINES_8514_8642_English.texfrag


% BEGIN INCLUDED SGA4_English_current_buildable_source_20260906/source/01/drafts/SGA4_I_LINES_8643_8705_English.texfrag
\hypertarget{sga4.exp.I.appendix.II.sec.6.lemma.2}{}
\begin{enonce*}{引理 2}
对每个基数 $\sheaf{\sheaf{c}}$, 都存在一个 $\sheaf{\sheaf{c}}$ 型且势为 $\sheaf{\sheaf{c}}$ 的 Artin 集 $B$.
\end{enonce*}

对 $\sheaf{\sheaf{c}}$ 作超限归纳. 当 $\sheaf{\sheaf{c}}=0$ 时别无选择, 取 $B=∅$. 若 $\sheaf{\sheaf{c}}$ 有前驱 $\sheaf{\sheaf{c}}'$, 设 $B'$ 为一个 $\sheaf{\sheaf{c}}'$ 型且势为 $\sheaf{\sheaf{c}}'$ 的 Artin 集; 此时 $\sheaf{\sheaf{c}}\leq2^{\sheaf{\sheaf{c}}'}$, 因而存在 $\mathcal{P}(B')$ 的势为 $\sheaf{\sheaf{c}}$ 的子集 $B$; $B$ 的元素是 $B'$ 的子集, 故其势为 $\leq\sheaf{\sheaf{c}}'\leq\sheaf{\sheaf{c}}$; $B$ 的更高阶分量是 $B'$ 的分量, 因而其势也为 $\leq\sheaf{\sheaf{c}}'\leq\sheaf{\sheaf{c}}$. 最后, 若 $\sheaf{\sheaf{c}}$ 没有前驱, 则对每个基数 $\sheaf{\sheaf{c}}_\lambda<\sheaf{\sheaf{c}}$, 选取一个 $\sheaf{\sheaf{c}}_\lambda$ 型且势为 $\sheaf{\sheaf{c}}_\lambda$ 的 Artin 集 $B_\lambda$; 此时 $B=\bigcup_\lambda B_\lambda$ 有所需性质. 这就证明了\hyperlink{sga4.exp.I.appendix.II.sec.6.lemma.2}{引理 2}, 并完成了 a) 的证明.

现在证明 b). 设 $\sheaf{\sheaf{c}}$ 为强不可达基数. 由 \hyperlink{sga4.exp.I.appendix.II.sec.6.theorem.2.part.a}{a)}, 已知严格 $\sheaf{\sheaf{c}}$ 型的 Artin 集构成集合 $U_{\sheaf{\sheaf{c}}}$. $U_{\sheaf{\sheaf{c}}}$ 为宇宙, 直接由 Artin 集的性质 \hyperlink{sga4.exp.I.appendix.II.sec.6.property.AR.I}{(AR.I)} 至 \hyperlink{sga4.exp.I.appendix.II.sec.6.property.AR.IV}{(AR.IV)} (\hyperlink{sga4.exp.I.appendix.II.sec.6}{本节}开头) 以及\pageoriginale 类似于\hyperlink{sga4.exp.I.appendix.II.sec.5}{第 5 节}\hyperlink{sga4.exp.I.appendix.II.sec.5.proposition.10}{命题 10} 中的基数估计得出. 对 $<\sheaf{\sheaf{c}}$ 的基数应用\hyperlink{sga4.exp.I.appendix.II.sec.6.lemma.2}{引理 2}, 就得到关系 $\sup_{x\in U_{\sheaf{\sheaf{c}}}}\card(x)=\sheaf{\sheaf{c}}$. 最后, 为证明 $\card(U_{\sheaf{\sheaf{c}}})=\sheaf{\sheaf{c}}$, 先假定 $\sheaf{\sheaf{c}}$ \emph{不可数}; 由\hyperlink{sga4.exp.I.appendix.II.sec.5.lemma.of.section.5}{第 5 节的引理}, $U_{\sheaf{\sheaf{c}}}$ 是并集 $\bigcup_{\sheaf{\sheaf{d}}<\sheaf{\sheaf{c}}}A_{\sheaf{\sheaf{d}}}$, 其中 $A_{\sheaf{\sheaf{d}}}$ 表示 $\sheaf{\sheaf{d}}$ 型的 Artin 集所成的集合; 而 $\card(A_{\sheaf{\sheaf{d}}})=2^{\sheaf{\sheaf{d}}}<\sheaf{\sheaf{c}}$ (由 \hyperlink{sga4.exp.I.appendix.II.sec.6.theorem.2.part.a}{a)}); 另一方面, 基数 $\sheaf{\sheaf{d}}<\sheaf{\sheaf{c}}$ 所成的集合的势为 $\leq\sheaf{\sheaf{c}}$ (第 III 章); 因而 $\card(U_{\sheaf{\sheaf{c}}})\leq\sheaf{\sheaf{c}}\cdot\sheaf{\sheaf{c}}=\sheaf{\sheaf{c}}$, 并且 $\card(U_{\sheaf{\sheaf{c}}})\geq\sheaf{\sheaf{c}}$, 因为对每个 $\sheaf{\sheaf{d}}<\sheaf{\sheaf{c}}$, 有 $\card(U_{\sheaf{\sheaf{c}}})\geq\card(A_{\sheaf{\sheaf{d}}})=2^{\sheaf{\sheaf{d}}}$.


$\sheaf{\sheaf{c}}=0$ 的情形显然, 因而只剩 $\sheaf{\sheaf{c}}$ 为\emph{可数无限基数}的情形. 在这种情形, $U_{\sheaf{\sheaf{c}}}$ 是有限型 (即自身及所有分量都有限) 的 Artin 集所成的集合, 此时用到一个令人愉快的组合结果.


% END INCLUDED SGA4_English_current_buildable_source_20260906/source/01/drafts/SGA4_I_LINES_8643_8705_English.texfrag


% BEGIN INCLUDED SGA4_English_current_buildable_source_20260906/source/01/drafts/SGA4_I_LINES_8706_8780_English.texfrag
\hypertarget{sga4.exp.I.appendix.II.sec.6.lemma.3}{}
\begin{enonce*}{引理 3}
{\rm (D. K\"onig?)}. 考虑由有限集 $E_n$ 和映射 $f_n:E_{n+1}\rightarrow E_n$ 组成的两个无限序列 $(E_n)_{n\geq1}$, $(f_n)_{n\geq1}$. 若不存在无限序列 $(x_n)_{n%%
\geq1}$ 使得 $x_n\in E_n$, 且对每个 $n$ 有 $f_n(x_{n+1})=x_n$, 则当 $n$ 充分大时 $E_n$ 为空.
\end{enonce*}

换言之, 若满足 $f_n(x_{n+1})=x_n$ 的一切序列 $(x_n)$ 都\emph{有限}, 则它们的长度\emph{有界}. 这可用投射极限表述: 非空有限集的投射极限非空 (参见 General Topology, 第 I 章, 第 2 版, \S 9, 第 6 节, 命题 8, $2^{\circ}$)).

事实上,\pageoriginale 用反证法. 若对任意大的 $n$, 都存在非空的 $E_n$, 则没有哪个 $E_n$ 为空 (因为存在 $f_n:E_{n+1}\rightarrow E_n$, 故 $E_n=∅$ 蕴涵 $E_{n+1}=∅$). 若有限序列 $(x_j)_{1\leq j\leq n}$ 满足对 $j=1,\ldots,n-1$ 有 $f_j(x_{j+1})=x_j$, 就称它\emph{相容}. 对 $n$ 作归纳, 证明存在相容序列 $(a_1,\ldots,a_n)$ $(a_i\in E_i)$, 使得对每个 $q\geq n$, 它都能扩充为长度 $q$ 的相容序列 $(a_1,\ldots,a_n,x_{n+1},\ldots,x_q)$. 当 $n=0$ 时这是显然的, 因为没有哪个 $E_q$ 为空. 从 $n$ 推至 $n+1$. 若对每个 $x\in f_n^{-1}(\{a_n\})\subset E_{n+1}$, 扩充 $(a_1,\ldots,a_n,x)$ 得到的所有相容序列的长度都以某整数 $q(x)$ 为界, 则由于 $E_{n+1}$ 有限, 扩充 $(a_1,\ldots,a_n)$ 得到的所有相容序列的长度也有界 (以 $\sup_xq(x)$ 为界); 因而存在 $a_{n+1}\in f_n^{-1}(\{a_n\})$, 使得相容序列 $(a_1,\ldots,a_n,a_{n+1})$ 可以扩充到任意长度. 这样就得到与假设矛盾的无限序列 $(a_n)_{n%%
\geq1}$.

由\hyperlink{sga4.exp.I.appendix.II.sec.6.lemma.3}{引理 3} 可知, 若 $A$ 为有限型的 Artin 集, 则它的链所成的有序集 $G(A)$ \emph{有限}: 事实上, 取具有 $n+1$ 项的链 $(x_n\in x_{n-1}\in\ldots\in x_1\in A)$ 所成的集合为 $E_n$ (因为 $A$ 的 $\leq n+1$ 阶分量的个数有限且每个都是有限集, 故此集有限), 并取映射 $(x_{n+1}\in x_n\in x_{n-1}\in\ldots)\mapsto(x_n\in x_{n-1}\in\ldots)$ 为 $f_n$. 而有限\pageoriginale 有序集的同构类集是\emph{可数}的: 事实上, 在 $\NN$ 的有限子集上给定序结构, 等价于给定它的图, 这是 $\NN\times\NN$ 的有限子集; 另一方面, 可数集的有限子集所成的集合可数 (第 III 章). 由 \hyperlink{sga4.exp.I.appendix.II.sec.6.theorem.2.proof.claim.star}{(*)} 及\hyperlink{sga4.exp.I.appendix.II.sec.6.lemma.1}{引理 1}, 可知有限型的 Artin 集所成的集合 $\mathfrak{a}$ 可数. 由\hyperlink{sga4.exp.I.appendix.II.sec.6.lemma.2}{引理 2} 可知它是无限集 (或者更简单地, 用 $z_0=∅$, $z_{n+1}=\{z_n\}$ 定义的序列 $(z_n)_{n\geq0}$). 这就完成了 b) 的证明.

\begin{enonce*}[remark]{注}
我们刚刚证明, 若 $A$ 为有限型的 Artin 集, 则它只有有限个分量. 因而存在整数 $n$, 使得 $A$ 为 $n$ 型.
\end{enonce*}
现在证明 c). 设 $\sheaf{\sheaf{c}}$ 为无限基数, $U$ 为含有势为 $\sheaf{\sheaf{c}}$ 的元素的宇宙, $A$ 为 $\sheaf{\sheaf{c}}$ 型的 Artin 集. 则有序集 $G(A)$ 的势为 $\leq\sheaf{\sheaf{c}}$ (上面的\hyperlink{sga4.exp.I.appendix.II.sec.6.theorem.2.proof.equation.2}{公式 (2)}). 将\hyperlink{sga4.exp.I.appendix.II.sec.6.lemma.1}{引理 1} 的第二个断言应用于 $G(A)$, 可知 $G(A)$ 上的映射 $(x_0,\ldots,x_n)\rightsquigarrow x_n$ 在 $U$ 中取值. 换言之, $A$ 的所有分量都是 $U$ 的元素. 这就证明了 c). \hyperlink{sga4.exp.I.appendix.II.sec.6.theorem.2.part.c-prime}{${\rm c'}$)} 的证明类似: 若 $A$ 为有限型的 Artin 集, 刚刚已经看到 $G(A)$ 有限; 因为 $U$ 非空, 由\hyperlink{sga4.exp.I.appendix.II.sec.1.corollary.to.proposition.7}{命题 7 的推论} (\hyperlink{sga4.exp.I.appendix.II.sec.1}{第 1 节}), 它含有与 $G(A)$ 等势的元素. \cqfd


% END INCLUDED SGA4_English_current_buildable_source_20260906/source/01/drafts/SGA4_I_LINES_8706_8780_English.texfrag


% BEGIN INCLUDED SGA4_English_current_buildable_source_20260906/source/01/drafts/SGA4_I_LINES_8781_8916_English.texfrag
\section{模糊的元数学说明}%
\hypertarget{sga4.exp.I.appendix.II.sec.7}{}\pageoriginale

\begin{enumerate}
\renewcommand{\theenumi}{\alph{enumi}}
\renewcommand{\labelenumi}{\theenumi)}
\item 公理「\emph{每个集合都是 Artin 集}」是\emph{无害的}. 事实上, 若有集合论的模型 $M$, 则 $M$ 的 Artin 元素所成的集合 $M'$ 也是模型 (参见\hyperlink{sga4.exp.I.appendix.II.sec.6.proposition.11}{命题 11}).

\item 宇宙公理 \hyperlink{sga4.exp.I.appendix.II.sec.4.axiom.A.6}{(A.6)} (以及关于强不可达基数的等价公理 \hyperlink{sga4.exp.I.appendix.II.sec.5.axiom.A-prime.6}{$(A'.6)$}) \emph{独立}于集合论的其余部分. 事实上, 设 $\sheaf{\sheaf{c}}$ 为第一个不可数强不可达基数. 严格 $\sheaf{\sheaf{c}}$ 型的 Artin 集所成的宇宙 $U_{\sheaf{\sheaf{c}}}$ (\hyperlink{sga4.exp.I.appendix.II.sec.6.theorem.2}{定理 2}, \hyperlink{sga4.exp.I.appendix.II.sec.6.theorem.2.part.b}{b)}) 是不含 \hyperlink{sga4.exp.I.appendix.II.sec.4.axiom.A.6}{(A.6)} 的集合论的模型: $U_{\sheaf{\sheaf{c}}}$ 的元素称为「集合」, 「隶属关系」是通常的隶属关系在 $U_{\sheaf{\sheaf{c}}}$ 上的限制, 等等. 因此模型的「宇宙」是作为 $U_{\sheaf{\sheaf{c}}}$ 的元素的通常意义的宇宙. 但我们已经看到, 作为 $U_{\sheaf{\sheaf{c}}}$ 的元素的宇宙, 只有两个小家伙 $U_0=∅$ 和 $U_1$. 因而 $U_1$ 是一个不属于任何「宇宙」的「集合」. 所以我们有了一个使 \hyperlink{sga4.exp.I.appendix.II.sec.4.axiom.A.6}{(A.6)} 为假的集合论模型.

\item Bourbaki 对\emph{无穷公理} (A.5) \emph{独立}于前面的公理与公理模式, 仅仅作了「推测」, 这太谨慎了. 它确实独立, 因为有限型的 Artin 集所成的可数宇宙 $U_1$ 是一个使 (A.5) 为假而前面的公理与公理模式为真的模型.

\item 若能证明宇宙公理 \hyperlink{sga4.exp.I.appendix.II.sec.4.axiom.A.6}{(A.6)} 是无害的, 将十分有趣. 这看来很难---其实 Paul Cohen 说, 这是不可证明的.
\end{enumerate}

标题中的形容词\pageoriginale「模糊」是说, 构造模型时, 我们没有费力摆弄符号 $\tau$, 使它不能将我们带到模型之外. 如果考虑一个模型 $M$, 诀窍是把通常的 $\tau_x(ℛ(x))$ 换成:
$$
\tau_x(ℛ(x)\mbox{ 且 } x\in M){\quoi}.
$$
还须检查, 这确实把通常的量词变为以前称为「典型」的量词:
$$
(\forall x \in M)\quav (\exists x \in M){\quoi}.
$$

\hypertarget{sga4.exp.I.appendix.II.sec.7.exercises}{}
\begin{enonce*}[remark]{练习}
\begin{enumerate}
\renewcommand{\labelenumi}{\theenumi)}
\hypertarget{sga4.exp.I.appendix.II.sec.7.exercise.1}{}
\item 设 $n$ 为整数 $\geq 1$. 证明 $n$ 型的 Artin 集所成的集合是无限集 (集合 $z_0=∅$, $z_1=\{∅\}$, $z_{q+1}=\{z_q\}$ 都是 1 型的).

\hypertarget{sga4.exp.I.appendix.II.sec.7.exercise.2}{}
\item 定义集合 $A$ 的\emph{高度}为使得存在序列 $x_n\in x_{n-1}\in \ldots \in x_0=A$ 的那些整数 $n$ 的上确界 (有限或无限). 证明高度 $\leq n$ 的集合都是有限集, 且构成一个有限集, 其势 $p_n$ 可由 $p_0=1$, $p_{n+1}=2^{p_n}$ 算出 (对 $n$ 作归纳, 注意高度 $\leq n$ 的集合 $A$ 的元素都是高度 $\leq n-1$ 的集合).

\hypertarget{sga4.exp.I.appendix.II.sec.7.exercise.3}{}
\item 设 $G$ 为具有最小元 $g_0$ 的 Noether 有序集, 并假定对每个 $h\in G$, 满足 $g\leq h$ 的 $g\in G$ 所成的集合有限且全序.
\begin{enumerate}
\renewcommand{\theenumii}{\alph{enumii}}
\renewcommand{\labelenumii}{\theenumii)}
\hypertarget{sga4.exp.I.appendix.II.sec.7.exercise.3.part.a}{}
\item 证明\pageoriginale $G$ 的每个元素 $h\neq g_0$ 都有唯一前驱.

\hypertarget{sga4.exp.I.appendix.II.sec.7.exercise.3.part.b}{}
\item 设 $\varphi$ 为\hyperlink{sga4.exp.I.appendix.II.sec.6.lemma.1}{引理 1} 中定义的 $G$ 上的映射 (即 $\varphi(g)$ 为 $g'$ 遍历 $g$ 的后继时, 一切 $\varphi(g')$ 所成的集合). 置 $A=\varphi(g_0)$. 证明 $A$ 为 Artin 集. 对 $g\in G$, 设 $g_0<g_1<\ldots<g_n=g$ 为 $\leq g$ 的元素所成的序列; 置 $f(g)=(\varphi(g_0),\varphi(g_1),\ldots,\varphi(g_n))$; 证明 $f$ 为从 $G$ 到正文中有序集 $G(A)$ 的递增满射. 证明若 $f$ 为\emph{单射}, 则它是从 $G$ 到 $G(A)$ 的\emph{同构}.

\hypertarget{sga4.exp.I.appendix.II.sec.7.exercise.3.part.c}{}
\item 对 $g\in G$, 设 $M_g$ 为 $g$ 的上界集. 假定对具有同一前驱 $p$ 的每对互异元素 $g$, $g'$, 有序集 $M_g$ 与 $M_{g'}$ 均不同构. 证明此时 \hyperlink{sga4.exp.I.appendix.II.sec.7.exercise.3.part.b}{b)} 的映射 $f$ 是从 $G$ 到 $G(A)$ 的\emph{同构}. [若 $f(g)=f(g')$ 且 $g\neq g'$, 并设 $g_0<g_1<\ldots<g_n=g$ 与 $g'_0<g'_1<\ldots <g'_{n'}=g'$ 分别为 $\leq g$ 及 $\leq g'$ 的元素所成的序列, 证明 $n=n'$, 并且可以假定 $g$ 与 $g'$ 具有同一前驱 $p$; 然后考虑满足下述条件的 $p\in G$ 所成的集合: $p$ 存在两个互异后继 $g$, $g'$, 使得 $f(g)=f(g')$; 考虑此集的一个极大元 $q$, 以及 $q$ 的两个互异后继 $h$, $h'$, 使得 $f(h)=f(h')$; 注意 $f$ 在 $M_h$ 与 $M_{h'}$ 上的限制均为单射, 因而 (由 \hyperlink{sga4.exp.I.appendix.II.sec.7.exercise.3.part.b}{b)}) 分别是从 $M_h$ 到 $G(\varphi(h))$, 及从 $M_{h'}$ 到 $G(\varphi(h'))$ 的同构; 由 $M_h$ 与 $M_{h'}$ 不同构的假设推出 $\varphi(h)\neq \varphi(h')$, 这与 $f(h)=f(h')$ 矛盾].
\end{enumerate}
注意.\pageoriginale \hyperlink{sga4.exp.I.appendix.II.sec.7.exercise.3}{练习 3}) 对 Artin 集如何「造出」给出了非常精确的信息. 我们已经知道, 这样的集合 $A$ 由有序集 $G(A)$ 的同构类确定 (\hyperlink{sga4.exp.I.appendix.II.sec.6.lemma.1}{引理 1}). 现在知道了如何\emph{刻画}与某个 $G(A)$ 同构的有序集 $G$:
\begin{itemize}
\item[$\alpha)$] $G$ 为 Noether 集且有最小元;

\item[$\beta)$] 对每个 $g\in G$, 满足 $h\leq g$ 的 $h\in G$ 所成的集合全序且有限 (从而每个非最小的 $g$ 的前驱都存在且唯一);

\item[$\gamma)$] 若 $g$, $g'$ 为具有同一前驱的互异元素, 则 $g$ 的上界集 $M_g$ 与 $g'$ 的上界集 $M_{g'}$ 不同构.
\end{itemize}

\hypertarget{sga4.exp.I.appendix.II.sec.7.exercise.4}{}
\item 设 $(A_n)_{n\geq 0}$ 为由 $A_0=∅$, $A_{n+1}=A_n\cup \{A_n\}$ 定义的有限\emph{序数}序列. 证明有序集 $G(A_n)$ 有 $2^n$ 个元素, 且其中高度为 $q$ 的元素 (\hyperlink{sga4.exp.I.appendix.II.sec.7.exercise.2}{练习 2} 的意义下) 的个数为 $\binom{n}{q}$.
\end{enumerate}
\end{enonce*}

% END INCLUDED SGA4_English_current_buildable_source_20260906/source/01/drafts/SGA4_I_LINES_8781_8916_English.texfrag


% BEGIN INCLUDED SGA4_English_current_buildable_source_20260906/source/01/drafts/SGA4_I_LINES_8917_8927_English.texfrag
\makeatletter
\let\toclevel@section\sgaFourSavedToclevelSection
\makeatother
\pdfbookmark[2]{参考文献}{sga4.exp.I.appendix.II.bibliography}
\begingroup
\makeatletter
\def\Hy@writebookmark#1#2#3#4#5{}
\makeatother
\begin{thebibliography}{10}%%
\bibitem[1]{sga4.exp.I.bib.I.2} \sisi{B. \textsc{Mitchell}, Theory of categories, Academic Press
(1965).}{B. \textsc{Mitchell}, \textit{Theory of categories},
Academic Press, 1965.}
\end{thebibliography}
\endgroup
\hypersetup{bookmarksdepth=1}

\ifx\danslelivre\undefined
%\bibliography{../fichiers-configuration-latex/bibliographie}
%\bibliographystyle{../fichiers-configuration-latex/smfalpha}
\end{document}
\fi

% END INCLUDED SGA4_English_current_buildable_source_20260906/source/01/drafts/SGA4_I_LINES_8917_8927_English.texfrag


